\documentclass[11pt,a4paper]{article}
\usepackage[utf8]{inputenc}
\usepackage[T1]{fontenc}
\usepackage{lmodern}
\usepackage{amsmath,amssymb,amsthm}
\usepackage{geometry}
\usepackage{enumitem}
\usepackage{fancyhdr}
\usepackage{titlesec}
\usepackage{array}
\usepackage{booktabs}
\usepackage{longtable}
\usepackage{graphicx}
\usepackage{float}

\geometry{margin=0.75in}

\pagestyle{fancy}
\fancyhf{}
\fancyhead[C]{\small\leftmark}
\fancyfoot[C]{\thepage}
\renewcommand{\headrulewidth}{0.4pt}

\newtheorem{theorem}{Theorem}[section]
\newtheorem{lemma}[theorem]{Lemma}
\newtheorem{proposition}[theorem]{Proposition}
\newtheorem{corollary}[theorem]{Corollary}
\theoremstyle{definition}
\newtheorem{definition}[theorem]{Definition}
\newtheorem{remark}[theorem]{Remark}

\DeclareMathOperator{\Norm}{Norm}

\titleformat{\section}{\large\bfseries}{\thesection.}{0.5em}{}
\titleformat{\subsection}{\normalsize\bfseries}{\thesubsection.}{0.5em}{}

\emergencystretch=3em
\tolerance=600
\begin{document}

\setcounter{page}{423}

%============================================================
% PAPER 412: A MEMOIR ON CUBIC SURFACES (continued)
%============================================================
\section*{[412] A Memoir on Cubic Surfaces. (\textit{continued})}
\markboth{A Memoir on Cubic Surfaces.}{Cayley --- Cubic Surfaces.}
\addcontentsline{toc}{section}{[412] A Memoir on Cubic Surfaces (continued)}

\subsection*{Section X = 12--B$_{\mbox{\scriptsize 3}}$--C$_{\mbox{\scriptsize 2}}$.\ Article Nos.\ 127 to 136.}

\noindent\textbf{127.} Writing
\[
(a,b,c,d\!\hbox{\hspace{-0.5em}}\raisebox{0.3ex}{$\scriptstyle\frown$}\!X,Y)^3 = d(f_1X-Y)(f_2X-Y)(f_3X-Y),
\]
the planes are
\[
\begin{array}{lll}
X=0, \quad [0] & \quad Z=0, \quad [7] & \quad W=0, \quad [8] \\
f_1X-Y=0, \quad [14] & \quad f_2X-Y=0, \quad [25] & \quad f_3Z-Y=0, \quad [36];
\end{array}
\]
and the lines are
\[
\begin{array}{lll}
X=0,\ Y=0, \quad (0) & \quad f_1X-Y=0,\ Z=0, \quad (1) & \quad f_2X-Y=0,\ Z=0, \quad (2) \\
f_3X-Y=0,\ Z=0, \quad (3) & \quad f_1Z-Y=0,\ W=0, \quad (4) & \quad f_2Z-Y=0,\ W=0, \quad (5) \\
f_3Z-Y=0,\ W=0, \quad (6). & &
\end{array}
\]

\noindent\textbf{128.} There is no facultative line; $p'=b'=0$, $t'=0$; and hence also $\beta'=0$.

\noindent\textbf{129.} \textit{Hessian surface.} The equation is
\[
X^2ZW(cX+dY) - 3X\{(ac-b^2,ad-bc,bd-c^2\!\hbox{\hspace{-0.5em}}\raisebox{0.3ex}{$\scriptstyle\frown$}\!X,Y)^2\} = 0,
\]
so that the Hessian breaks up into the plane $X=0$ (axial or common biplane) and into a cubic surface.

The complete intersection of the Hessian with the cubic surface is made up of the line $X=0$, $Y=0$ (the axis) four times; and of a system of four conics, which is the spinode curve; $c'=8$.

In fact combining the equations $WXZ+(a,b,c,d\!\hbox{\hspace{-0.5em}}\raisebox{0.3ex}{$\scriptstyle\frown$}\!X,Y)^3=0$ and
\[
ZW(cX+dY)-3X(ac-b^2,ad-bc,bd-c^2\!\hbox{\hspace{-0.5em}}\raisebox{0.3ex}{$\scriptstyle\frown$}\!X,Y)^2=0,
\]
these intersect in the axis once, and in a curve of the eighth order which breaks up into four conics; for we can from the two equations deduce
\[
(a,b,c,d\!\hbox{\hspace{-0.5em}}\raisebox{0.3ex}{$\scriptstyle\frown$}\!X,Y)^3(cX+dY)+3X^2(ac-b^2,ad-bc,bd-c^2\!\hbox{\hspace{-0.5em}}\raisebox{0.3ex}{$\scriptstyle\frown$}\!X,Y)^2=0,
\]
that is $(4ac-3b^2,ad,bd,cd,d^2\!\hbox{\hspace{-0.5em}}\raisebox{0.3ex}{$\scriptstyle\frown$}\!X,Y)^4=0$, a system of four planes each intersecting the cubic $XZW+(a,b,c,d\!\hbox{\hspace{-0.5em}}\raisebox{0.3ex}{$\scriptstyle\frown$}\!X,Y)^3=0$ in the axis and a conic; whence, as above, spinode curve is four conics.

It is easy to see that the tangent planes along any conic on the surface pass through a point, and form therefore a quadric cone; hence in particular the spinode torse is made up of the quadric cones which touch the surface along the four conics respectively.

\setcounter{page}{424}

\noindent\textbf{130.} \textit{Reciprocal Surface.} The equation is obtained by means of the binary cubic $\lambda(\lambda x+yY)^2 + 12zw(a,b,c,d\!\hbox{\hspace{-0.5em}}\raisebox{0.3ex}{$\scriptstyle\frown$}\!X,Y)^3$, viz.\ calling this $({*}\!\hbox{\hspace{-0.5em}}\raisebox{0.3ex}{$\scriptstyle\frown$}\!X,Y)^3$ the coefficients are $(x^2+12azw,\ 2xy+12bzw,\ y^2+12czw,\ 12dzw)$. The equation is found to be
\begin{multline*}
-432(d^2a^2-6abcd+4ac^3+4b^3d-3b^2c^2)z^3w^3 \\
+216\{(ad^2-3bcd+2c^3)x^2+(-2acd+4b^2d-2bc^2)xy+(-abd+2ac^2-b^2c)y^2\}zw \\
+9\{(d^2a^2-12cdaxy+(10bd+3c^2)axy^2-(4ad+8bc)xy^3+(4ac-b^2)y^4\}zw \\
-y^3(da^3-3caxy+3bxy^2-ay^3)=0.
\end{multline*}
The section by the plane $w=0$ (reciprocal of $\mathrm{B}_1=\mathrm{D}$) is the line $w=0$, $y=0$ (reciprocal of edge) three times, and the lines $w=0$, $dx^2-3caxy+3bxy^2-ay^3=0$ (reciprocals of the biplanar rays). And similarly for the section by the plane $z=0$ (reciprocal of $\mathrm{B}_2=\mathrm{C}$).

The section by the plane $y=0$ is made up of the lines $(y=0,\ z=0)$, $(y=0,\ w=0)$ each once, and of two conics, $y=0$,
\[
16(d^2a^2-6abcd+4ac^3+4b^3d-3b^2c^2)z^2w^2 + 8(ad^2-3bcd+2c^3)x^2zw + x^4 = 0.
\]

\noindent\textbf{131.} There is not any nodal curve; $b'=0$.

\noindent\textbf{132.} \textit{Cuspidal curve.} The equations may be written
\[
\begin{vmatrix}
x^2+12azw & 2xy+12bzw & y^2+12czw \\
2xy+12bzw & y^2+12czw & 12dzw
\end{vmatrix}=0.
\]
Forming the equations
\begin{align*}
(bd-c^2)\cdot 12zw + 2(dxy-cy^2)\cdot 12zw - y^4 &= 0,\\
(ad-bc)\cdot 12zw + (3dx^2-2cxy-by^2)\cdot 12zw - 2xy^2 &= 0,
\end{align*}
these are two quartic surfaces having in common the lines $(y=0,\ w=0)$, $(y=0,\ z=0)$: attending to the line $(y=0,\ z=0)$, this is on the second surface an oscular line, $z=\frac{1}{2}y$; on the first surface it is a nodal line, the one tangent plane being $6(bd-c^2)w\cdot z+dxz\cdot y=0$, the other tangent plane being $z=0$, but the line being in regard to this sheet an oscular line, $z=\frac{1}{2}y-\frac{y^2}{2}$. Hence in the intersection of the two surfaces the line counts $(1+3=)4$ times; similarly the line $y=0$, $w=0$ counts $(1+3=)$ four times; and there is a residual intersection of the order $(16-4-4=)8$, which is the cuspidal curve; $c'=8$.

\setcounter{page}{425}

\noindent\textbf{133.} The cuspidal curve is a system of four conics; in fact from the preceding equations written in the forms
\begin{align*}
(bd-c^2,\ 2(dxy-cy^2),\ -y^4\!\hbox{\hspace{-0.5em}}\raisebox{0.3ex}{$\scriptstyle\frown$}\!12zw,\ 1)^2 &=0,\\
(ad-bc,\ 3dx^2-2cxy-by^2,\ -2xy^2\!\hbox{\hspace{-0.5em}}\raisebox{0.3ex}{$\scriptstyle\frown$}\!12zw,\ 1)^2 &=0,
\end{align*}
eliminating $zw$, we obtain
\[
\{4^2(bd-c^2),\ -2(ad^2-3bcd+4c^3),\ 6(acd+3b^2d-2bc^2),\ 6(bc-ad)b,\ a^2d-3abc+2b^3\!\hbox{\hspace{-0.5em}}\raisebox{0.3ex}{$\scriptstyle\frown$}\!x,y\}^4=0,
\]
which shows that the cuspidal curve lies in four planes, and it hence consists of four conics; these are of course the reciprocals of the quadric cones which touch the cubic surface along the four conics which make up the spinode curve.

\noindent\textbf{134.} The equation of the surface, attending only to the terms of the second order in $y$, $z$, $w$, is $27d^2z^2w^2=0$; it thus appears that the point $y=0$, $z=0$, $w=0$ (reciprocal of the plane $X=0$) (which is oscular along the axis joining the two binodes, or $\mathrm{B}_5$-axis) is a binode on the reciprocal surface, the biplanes being $z=0$, $w=0$, viz.\ these are the planes reciprocal to the binodes $(Z=0,\ Y=0,\ W=0)$ and $(Z=0,\ Y=0,\ Z=0)$ of the cubic surface; we have thus $\mathrm{B}'=1$.

It is proper to remark that the binode $y=0$, $z=0$, $w=0$ is not on the cuspidal curve, as its being so would probably imply a higher singularity.

\noindent\textbf{135.} A simple case, presenting the same singularities as the general one, is when $a=d$, $b=c=0$: to diminish the numerical coefficients assume $a=d=\frac{1}{2}$, the cubic surface is thus $12XZW+X^3+Y^3=0$, and the equation of the reciprocal surface, multiplying it by 4, becomes $z^2w^2+6x^2zw^2+(3x^4-12x^2y^2)zw-4y^2(x^4-y^4)=0$, viz.\ this is the surface $4y^6-4y^2x^2(x^2+6zw)+zw(3x^2+zw)^2=0$ considered in the Memoir ``On the Theory of Reciprocal Surfaces.''

The cuspidal curve is, as there shown, composed of the four conics $y=0$, $3x^2+4zw=0$ and $y^2-4x^2=0$, $x^2-4zw=0$; and it is there shown that the two points $(x=0,\ y=0,\ z=0)$, $(x=0,\ y=0,\ w=0)$, each reckoned eight times, are to be considered as off-points of the reciprocal surface.

\setcounter{page}{426}

\noindent\textbf{136.} The like investigation applies to the general surface, and we have thus $\chi'=16$; the points in question are still the points $(x=0,y=0,z=0)$, $(x=0,y=0,w=0)$; viz.\ these are the points of intersection of the surface by the line $(x=0,\ y=0)$, which points are also the common points of intersection of the four conics which compose the cuspidal curve, that is, they are quadruple points on the cuspidal curve; it does not appear that the points are on this account, viz.\ as quadruple points of the cuspidal curve, off-points of the surface, nor does this even show that the points should be reckoned each eight times. As already remarked, the singularity requires a more complete investigation.

\subsection*{Section XI = 12--B$_{\mbox{\scriptsize 3}}$.\ Article Nos.\ 137 to 143.\ Equation $WXZ+(X+Z)(Y^2-X^2)=0$.}

\noindent\textbf{137.} The diagram of the lines and planes is shown in the original text. The arrangement consists of:
\begin{itemize}[noitemsep]
\item Biplane touching along axis and containing edge ($1\times 12=12$)
\item Other biplane ($1\times 12=12$)
\item Planes each through the axis and containing a biplanar ray and a cnienodal ray ($2\times 8=16$)
\item Plane touching along the edge and containing the mere line ($1\times 3=3$)
\item Biradial plane through the two cnienodal rays ($1\times 2=2$)
\end{itemize}
Total: 45 planes.

\setcounter{page}{427}

\noindent\textbf{138.} The planes are
\[
\begin{array}{ll}
X=0, \quad [0] & Z=0,\ Y=0, \quad (0) \\
Z=0, \quad [3] & X=0,\ Z=0, \quad (3) \\
X-Y=0, \quad [11'] & X-Y=0,\ Z=0, \quad (1) \\
Z+Y=0, \quad [22'] & X+Y=0,\ Z=0, \quad (2) \\
W=0, \quad [31] & Z-Y=0,\ W=0, \quad (10) \\
Z+X=0, \quad [1'2'] & Z+Y=0,\ W=0, \quad (2') \\
& X+Z=0,\ W=0, \quad (12).
\end{array}
\]

\noindent\textbf{139.} The facultative lines are the edge counting twice, and the mere line; $p'=b'=3$; $t'=1$.

\noindent\textbf{140.} \textit{Hessian surface.} The equation is $X(X+Z)(ZW+3X^2-XZ)+Y^2(Z-X)=0$. The complete intersection with the surface consists of the line $(X=0,\ Y=0)$, the axis, four times; the line $(X=0,\ Z=0)$, the edge, twice; and a sextic curve, which is the spinode curve; $c'=6$.

Writing the equations of the surface and the Hessian in the form
\begin{align*}
X(ZW+Y^2)-Z^3+Z(Y^2-X^2) &= 0,\\
X(X+Z)(ZW+Y^2)+(Z-3X)\{-X^2+Z(Y^2-X^2)\} &= 0,
\end{align*}
we see that the equations of the spinode curve may be written $ZW+Y^2=0$, $-X^2+Z(Y^2-X^2)=0$, viz.\ the curve is a complete intersection, $2\times 3$.

There is at $E_3$ a triple point, the tangents coinciding with the nodal rays $W=0$, $Y^2-X^2=0$. The edge and the mere line are each of them single tangents of the spinode curve. But the edge counting twice in the nodal curve, its contact with the spinode curve will also count twice, that is, we have $\beta'=2.1+1=3$.

\setcounter{page}{428}

\noindent\textbf{141.} \textit{Reciprocal Surface.} The equation is obtained by means of the binary cubic $4w^2Z(X+Z)^2 + 4wZ(X+Z)(xX+zZ) + fXZ^2$, or calling this $({*}\!\hbox{\hspace{-0.5em}}\raisebox{0.3ex}{$\scriptstyle\frown$}\!X,Z)^2$, the coefficients are $(12w^2,\ 8w^2+4wx,\ 4w^2+4wx+4wz+y^2,\ 12wz)$, and thence the equation is found to be
\begin{multline*}
w^3[y^4-4(x-z)^2] + 16w^2[(2x-3z)y^2-2(x-z)(x-z)^2] \\
+8w[y^4+(3x^2-xz+6z^2)y^2-2x^2(x-z)^2] \\
+4w[(2x+3z)y^4-2x^2(x+z)y^2]+y^4(y^2-x^2)=0,
\end{multline*}
where the section by the plane $w=0$ (reciprocal of binode) is $y^4(y^2-x^2)=0$, viz.\ this is the line $w=0$, $y=0$ (reciprocal of the edge) four times, and the lines $w=0$, $y^2-x^2=0$ (reciprocals of the biplanar rays).

The section by the plane $z=0$ is found to be $(y^2-x^2)\{y^2+4xw+4w^2\}^2=0$, viz.\ this is the two lines $z=0$, $y^2-x^2=0$ (reciprocals of the nodal rays), and the conic $z=0$, $y^2+4xw+4w^2=0$ (reciprocal of the nodal cone $WX+Y^2-X^2=0$) twice.

\noindent\textbf{142.} \textit{Nodal curve.} The equation shows that the line $y=0$, $x-z=0$ (reciprocal of the line $W=0$, $X+Z=0$) is a nodal line on the surface. It also shows that the line $y=0$, $w=0$ (reciprocal of the edge) is a tacnodal line ($=2$ nodal lines) on the surface; in fact attending only to the lowest terms in $y$, $w$, we have $-4\{16(x-z)^2w^2+8(x+z)wy^2+y^4\}=0$, that is,
\[
y^2\left\{\frac{1}{4}y^2+2(x+z)w+16(x-z)^2\frac{w^2}{y^2}\right\}=0,
\]
two values, $w=Ay^2$, $w=By^2$, which indicates a tacnodal line. The nodal curve is thus made up of the line $y=0$, $x-z=0$ once, and the line $y=0$, $w=0$ twice; $b'=3$.

\noindent\textbf{143.} \textit{Cuspidal curve.} The equations give
\begin{align*}
(4w+2x)^2-3(4w^2+4wx+4wz+y^2) &=0,\\
-4w^2z+(2w+x)(4w^2+4wx+4wz+y^2) &=0,
\end{align*}
or, as these are more simply written,
\begin{align*}
w^2+4wx-12wz+4x^2-3y^2 &=0,\\
8w^3+12w^2x-28w^2z+w(4x^2+4xz+2y^2)+xy^2 &=0,
\end{align*}
so that the cuspidal curve is a complete intersection $2\times 3$; $c'=6$.


\subsection*{Section XII = 12--U$_{\mbox{\scriptsize 3}}$.\ Article Nos.\ 150 to 156.\ Equation $W(X+Y+Z)^2+XYZ=0$.}

\noindent\textbf{150.} The diagram of the lines and planes is: Lines include $3\times 1=3$ mere lines, and $1'2'3'$; Planes include the uniplane $[0]$, $1\times 32=32$; $3\times 4=12$ planes each touching along a ray and containing a mere line; and $1\times 1=1$ plane through the three mere lines. Total: 46 lines and planes.

\noindent\textbf{151.} The planes are $X+Y+Z=0$, $[0]$; $X=0$, $[1]$; $Y=0$, $[2]$; $Z=0$, $[3]$; $W=0$, $[1'2'3']$. The lines are $X=0$, $Y+Z=0$, $(1)$; $Y=0$, $Z+X=0$, $(2)$; $Z=0$, $X+Y=0$, $(3)$; $X=0$, $W=0$, $(1')$; $Y=0$, $W=0$, $(2')$; $Z=0$, $W=0$, $(3')$; and $W=0$, $X+Y+Z=0$, $(012)$.

\noindent\textbf{152.} The three mere lines are each facultative: $p'=b'=3$; $t'=1$.

\noindent\textbf{153.} \textit{Hessian surface.} The equation is $(X+Y+Z)^2(X^2+Y^2+Z^2-2YZ-2ZX-2XY)=0$, viz.\ the surface consists of the uniplane $X+Y+Z=0$ twice, and of a quadric cone having its vertex at $U_3$ and touching each of the planes $X=0$, $Y=0$, $Z=0$.

The complete intersection with the cubic surface is made up of the rays each twice and of a residual sextic, which is the spinode curve; $c'=6$. The equations of the spinode curve are
\[
W(X+Y+Z)^2+XYZ=0, \quad X^2+Y^2+Z^2-2YZ-2ZX-2XY=0,
\]
viz.\ the curve is a complete intersection, $2\times 3$. Each of the mere lines is a single tangent; that is, $\beta'=3$.

\setcounter{page}{432}

\noindent\textbf{154.} \textit{Reciprocal Surface.} The equation is found by means of the binary cubic $4(T-xU)(T-yU)(T-zU)+wT^2U$, viz.\ writing for shortness $\beta=x+y+z$, $\gamma=yz+zx+xy$, $\delta=xyz$, then the cubic function is $(12,\ w-4\beta,\ 4\gamma,\ -12\delta\!\hbox{\hspace{-0.5em}}\raisebox{0.3ex}{$\scriptstyle\frown$}\!T,U)^3$, and the equation of the reciprocal surface is found to be
\[
432\beta\delta^2 + 64\gamma^3 + 72(w-4\beta)\gamma\delta - (w-4\beta)^2\gamma^2=0;
\]
expanding, this is
\[
w^2(-\beta\delta) + 8w(-6\beta^2\delta+\beta\gamma^2+9\gamma\delta) + 16(4\beta^2\delta-\beta^2\gamma^2-18\beta\gamma\delta+4\gamma^3+27\delta^2)=0;
\]
or substituting for $\beta$, $\gamma$, $\delta$ in the first and last lines, this is
\begin{multline*}
w^2(-xyz) + w^2(12\beta\delta-\gamma^2) + 8w(-6\beta^2\delta+\beta\gamma^2+9\gamma\delta) \\
+ 16(y-z)^2(z-x)^2(x-y)^2=0
\end{multline*}
(where $\beta$, $\gamma$, $\delta=x+y+z$, $yz+zx+xy$, $xyz$).

The section by the plane $w=0$ (reciprocal of the unode) is made up of the lines $w=0$, $y-z=0$; $w=0$, $z-x=0$; $w=0$, $x-y=0$ (reciprocals of the rays) each twice.

\noindent\textbf{155.} The nodal curve is at once seen to consist of the lines $(y=0,\ z=0)$, $(z=0,\ x=0)$, $(x=0,\ y=0)$, reciprocals of the facultative lines; in fact, in regard to $(y,z)$ conjointly $\gamma$ is of the order 1, and $\delta$ is of the order 2; hence every term of the equation is of the order 2 in $y$, $z$; and the like as to the other two lines: $b'=3$ as above.

\setcounter{page}{433}

\noindent\textbf{156.} For the cuspidal curve we have
\[
\begin{vmatrix}
12 & w-4\beta & w-4\beta & 4\gamma \\
w-4\beta & 4\gamma & -12\delta
\end{vmatrix}=0,
\]
or say $48\gamma^2-(w-4\beta)^2=0$, $36\delta+\gamma(w-4\beta)=0$, whence the cuspidal curve is a complete intersection $2\times 3$; $c'=6$.

\subsection*{Section XIII = 12--B$_{\mbox{\scriptsize 3}}$--2C$_{\mbox{\scriptsize 2}}$.\ Article Nos.\ 157 to 164.\ Equation $WXZ+Y^2(Y+X+Z)=0$.}

\noindent\textbf{157.} The diagram of the lines and planes is: Lines include $1$ transversal, $2\times 6=12$; Planes include biplanes ($2\times 6=12$), a plane through the three axes ($1\times 4=4$), planes each through an axis joining the binode with a onicnode ($2\times 6=12$), planes through the biplanar rays ($1\times 3=3$), and a plane touching along the axis which joins the two onicnodes ($1\times 2=2$). Total: 45.

\noindent\textbf{158.} The planes are $X=0$, $[1]$; $Z=0$, $[2]$; $Y=0$, $[056]$; $Y+Z=0$, $[5]$; $Y+X=0$, $[6]$; $Y-W=0$, $[34]$; $X+Y+Z=0$, $[12]$; $W=0$, $[0]$; and $W=0$, $X+Y+Z=0$, $(012)$. The lines are $Z=0$, $Y=0$, $(5)$; $Z=0$, $X=0$, $(6)$; $Y=0$, $W=0$, $(0)$; $X=0$, $Y+Z=0$, $(1)$; $Z=0$, $Y+X=0$, $(2)$; $W=Y=-Z$, $(3)$; $W=Y=-X$, $(4)$.

\noindent\textbf{159.} The transversal is facultative; $p'=b'=1$, $t'=0$.

\noindent\textbf{160.} The Hessian surface is $WXZ(5Y+X+Z)+Y^2(Z-X)^2=0$. The complete intersection with the surface is made up of the line $Y=0$, $X=0$ (C$\xi$-axis) three times; the line $Y=0$, $Z=0$ (C$\eta$-axis) three times; line $Y=0$, $W=0$ (C$\zeta$-axis) twice, and of a residual quartic, which is the spinode curve; $c'=4$.

\noindent\textbf{161.} Representing the two equations by $U=0$, $H=0$, we have $(3Y+X+Z)U-H=Y\{3Y^2+4YZ+4XZ+4X^2\}$, and further reductions lead to the spinode curve as a complete intersection $2\times 2$; and since the first surface is a cone having its vertex on the second surface, the spinode curve is a nodal quadriquadric. The equations of the transversal are $W=0$, $X+Y+Z=0$, and substituting in the equations of the spinode curve we obtain from each equation $(X-Z)^2=0$, that is, the transversal is a single tangent of the spinode curve; $\beta'=1$.

\setcounter{page}{435}

\noindent\textbf{162.} \textit{Reciprocal Surface.} The equation is derived from that belonging to $\mathrm{X}=12-\mathrm{B}_3-\mathrm{C}_2$ by writing therein $a=b=0$, $c=\frac{1}{2}$, $d=1$. The equation becomes
\begin{multline*}
y^2\cdot 16(x-z)^2 + wy^2\{-y^2(x+z)\\
+2y(x^2-4xz+z^2)+(x+z)(2x-z)(-x+2z)\} \\
+w\{y^4+8y^3(x+z)-2y^2(3x^2+28xz+4z^2)\\
+36xyz(x+z)-27xz^2\}+y^4(y-x)(y-z)=0.
\end{multline*}
The section by the plane $w=0$ (reciprocal of B$_3$) is $w=0$, $y=0$ (the edge) three times; and $w=0$, $y-x=0$; $w=0$, $y-z=0$ (reciprocals of the CB-axes).

\noindent\textbf{163.} \textit{Nodal curve.} This is the line $y=x=z$; wherefore $b'=1$. To put the line in evidence, write for a moment $x=y+\alpha$, $z=y+\gamma$, then the equation is readily converted into a form which, each term being at least of the second order in $\alpha$, $\gamma$ ($=x-y$, $z-y$), exhibits the nodal line in question.

\setcounter{page}{436}

\noindent\textbf{164.} \textit{Cuspidal curve.} Multiplying by 27, the equation may be written
\[
(7y-3x-3z-8w,\ -y+6w,\ -4w\!\hbox{\hspace{-0.5em}}\raisebox{0.3ex}{$\scriptstyle\frown$}\!y^2+16yw-12(x+z)w+16w^2,\\
-20y^2+24y(x+z)-27xz-8yw+16w^2)=0,
\]
where $4w(7y-3x-3z-8w)+(-y+6w)^2=y^2+16yw-12(x+z)w+16w^2$; and we have thus in evidence the cuspidal curve,
\begin{align*}
y^2+16yw-12(x+z)w+16w^2 &= 0,\\
-20y^2+24y(x+z)-27xz-8yw+16w^2 &= 0,
\end{align*}
which is a complete intersection, $2\times 2$, or quadriquadric curve; $c'=4$.

\subsection*{Section XIV = 12--B$_{\mbox{\scriptsize 5}}$--C$_{\mbox{\scriptsize 2}}$.\ Article Nos.\ 165 to 171.\ Equation $WXZ+Y^2Z+YX^2=0$.}

\noindent\textbf{165.} The diagram of the lines and planes is: Planes include $Z=0$ (torsal biplane, $1\times 15=15$), $X=0$ (ordinary biplane, $1\times 20=20$), $Y=0$ (plane through axis and two rays, $1\times 10=10$). Total: 45.

\noindent\textbf{166.} The edge is a facultative line, as will appear from the discussion of the reciprocal surface: $p'=b'=1$; $t'=0$.

\noindent\textbf{167.} \textit{Hessian surface.} The equation is $WXZ^2+Y^2Z^2-3X^2YZ+X^4=0$. The complete intersection with the surface is made up of the line $X=0$, $Y=0$ (the axis) five times, the line $X=0$, $Z=0$ (the edge) four times, and a skew cubic, the equations of which may be written
\[
\begin{vmatrix}
X, & Y, & W \\
4Z, & X, & -5Y
\end{vmatrix}=0.
\]

\setcounter{page}{438}

\noindent\textbf{168.} The spinode curve is made up of the edge $X=0$, $Z=0$ once, and of the cubic curve; and therefore $c'=4$.

\noindent\textbf{169.} \textit{Reciprocal Surface.} The equation is found to be
\begin{multline*}
(y^2+4xw)^2(y+w-x-z) - 27z(y+2w)^3 \\
+ 9xz(y^2+4xw+4zw)(y+2w) - 27xz^2w = 0,
\end{multline*}
reducing to
\begin{multline*}
y^2\cdot 8(x-z)^2 + wy^2\{-y^2(x+z)+2y(x^2-4xz+z^2)\\
+(x+z)(2x-z)(-x+2z)\} + w\{y^4+8y^3(x+z)\\
-2y^2(3x^2+28xz+4z^2)+36xyz(x+z)-27xz^2\}\\
+y^4(y-x)(y-z)=0.
\end{multline*}
The section by the plane $w=0$ (reciprocal of B$_5$) is $w=0$, $y=0$ (reciprocal of edge) four times, and the lines $w=0$, $y^2+4xz=0$ (reciprocals of the rays).

\setcounter{page}{439}

\noindent\textbf{170.} \textit{Nodal curve.} The equation gives $16w^2$ showing that the line $w=0$, $y=0$ (reciprocal of edge) is an oscnodal line counting as three lines; $b'=3$. There is no cuspidal curve; $c'=0$.

\subsection*{Section XV = 12--U$_{\mbox{\scriptsize 5}}$.\ Article Nos.\ 172 to 177.\ Equation $X^2W-(XZ+Y^2)^2=0$.}

\noindent\textbf{172.} The diagram of the lines and planes is: Plane is $X=0$ ($1\times 45=45$), a uniplane.

\noindent\textbf{173.} There is no facultative line; $p'=b'=0$, $t'=0$.

\noindent\textbf{174.} The Hessian surface is $X^3Y^2=0$, viz.\ this is the uniplane $X=0$, three times, and the plane $Y=0$ through the ray. The complete intersection with the cubic surface is made up of $X=0$, $Y=0$ (the ray) ten times, and of a residual conic, which is the spinode curve; $c'=2$. The equations of the spinode conic are $Y=0$, $XW+Z^2=0$, viz.\ the plane of the conic passes through the ray. Since there is no facultative line, $\beta'=0$.

\setcounter{page}{440}

\noindent\textbf{175.} \textit{Reciprocal Surface.} The equation is at once found to be $27(z^2+4xw)^2y^2-64x^3w^3=0$. The section by the plane $w=0$ (reciprocal of the Unode) is $w=0$, $y=0$ (reciprocal of ray) four times. There is no nodal curve; $b'=0$. But there is a cuspidal conic $y=0$, $z^2+4xw=0$. The point $y=0$, $z=0$, $w=0$ (reciprocal of the uniplane $X=0$) is a point which must be considered as uniting the singularities $\mathrm{B}'=1$, $\chi'=2$.

\subsection*{Section XVI = 12--4C$_{\mbox{\scriptsize 2}}$.\ Article Nos.\ 178 to 183.\ Equation $X^2W+XYZ+Y^3+Z^3=0$.}

\noindent\textbf{178.} There is no facultative line; $p'=b'=0$, $t'=0$.

\noindent\textbf{179.} \textit{Hessian surface.} The equation is $Z(X^2W+XYZ+Y^3-Z^3)=0$, viz.\ this breaks up into the plane $Z=0$ (the common biplane) and a cubic surface. The complete intersection with the surface is made up of the axes each four times, and of a residual sextic, which is the spinode curve; $c'=6$.

\noindent\textbf{180.} \textit{Reciprocal Surface.} The equation is found to be
\begin{multline*}
(x^2-4yz)^2\{x^3+4w(y^2+z^2)\} + 54w^2(x^2-4yz)(x^2+2yz) \\
+ 27w^3(27w-8x^3-24xyz)=0.
\end{multline*}
The section by the plane $w=0$ (reciprocal of C$_2$) is $w=0$, $x^2-4yz=0$ (reciprocal of nodal cone) four times. There is no nodal curve; $b'=0$.

\setcounter{page}{441}

\noindent\textbf{181.} \textit{Cuspidal curve.} The cuspidal curve is a complete intersection $2\times 3$; $c'=6$.

\subsection*{Section XVII = 12--2B$_{\mbox{\scriptsize 3}}$--C$_{\mbox{\scriptsize 2}}$.\ Article Nos.\ 184 to 189.}

\noindent\textbf{184--189.} In this case there is no facultative line; $p'=b'=0$, $t'=0$. The Hessian surface breaks up into the common biplane and a cubic surface. The spinode curve is a conic; $c'=2$. The reciprocal surface has no nodal curve but has a cuspidal conic; $c'=2$.

\subsection*{Section XVIII = 12--B$_{\mbox{\scriptsize 4}}$--2C$_{\mbox{\scriptsize 2}}$.\ Article Nos.\ 190 to 193.}

\noindent\textbf{190--193.} The edge counting twice is facultative; $p'=b'=3$, $t'=1$. The reciprocal of the edge is a tacnodal line. The spinode curve and cuspidal curve are both nonexistent; $c'=0$.

\subsection*{Section XIX = 12--B$_{\mbox{\scriptsize 5}}$--C$_{\mbox{\scriptsize 3}}$.\ Article Nos.\ 194 to 197.}

\noindent\textbf{194--197.} The axis counting three times is facultative; $p'=b'=3$, $t'=0$. The reciprocal of the axis is oscnodal. There is no spinode curve; $c'=0$.

\subsection*{Section XX = 12--U$_{\mbox{\scriptsize 5}}$.\ Article Nos.\ 198 to 201.}

\noindent\textbf{198--201.} There is no facultative line; $p'=b'=0$, $t'=0$. The Hessian surface is the uniplane three times, and a plane through the ray. The spinode curve is a conic; $c'=2$. The reciprocal surface has no nodal curve, but has a cuspidal conic.

\subsection*{Section XXI = 12--3B$_{\mbox{\scriptsize 3}}$.\ Article Nos.\ 202 to 204.}

\noindent\textbf{202--204.} There is no facultative line; $p'=b'=0$, $t'=0$. The Hessian surface consists of the common biplane and the other biplanes each once. The complete intersection consists of the axes each four times; there is no spinode curve, $c'=0$; whence also $\beta'=0$.

\setcounter{page}{451}

\subsection*{Article No.\ 202.\ Synopsis for the foregoing sections.}

\noindent\textbf{202.} I annex the following synopsis, for the several cases, of the facultative lines (or node-couple curve) and of the spinode curve of the cubic surface; also of the nodal curve and the cuspidal curve of the reciprocal surface.

It is to be observed that in designating a curve, for instance, as $18=4\times 5-2$, this means that it is a curve of the order 18, the partial intersection of a quartic surface and a quintic surface, but without any explanation of the nature of the common curve 2 which causes the reduction, viz.\ without explaining whether this is a conic or a pair of lines, and so in other cases; this may be seen by reference to the proper section of the Memoir.

\begin{center}
\begin{tabular}{lllll}
\toprule
& \textit{Facultative lines.} & \textit{Nodal curve.} & \textit{Spinode curve.} & \textit{Cuspidal curve.} \\
\midrule
I = 12 & 27 & 27 & $12=3\times 4$ & $24=6\times 4$ \\
II = 12--C$_2$ & 15 & 16 & $12=3\times 4$ & $18=4\times 5-2$ \\
III = 12--B$_3$ & 9 & 9 & $12=3\times 4$ & $16=4\times 5-4$ \\
IV = 12--2C$_2$ & 7 & 7 & $10=3\times 4-2$ & $12=4\times 4-2-2$ \\
V = 12--B$_4$ & $7=5+\text{edge twice}$ & $7=5+\text{rec.\ of edge twice}$ & $10=3\times 4-2$ & $12=4\times 4-4$ \\
VI = 12--B$_3$--C$_2$ & 3 & 3 & $9=3\times 4-3$ & $10=4\times 4-4-2$ \\
VII = 12--B$_5$ & $3=2+\text{edge}$ & $3=2+\text{rec.\ of edge}$ & $9=\text{edge}+\text{unicursal 8-thic}$ & $10=\text{rec.\ of edge}+\text{unicursal 9-thic}$ \\
VIII = 12--3C$_2$ & 3 & 8 & $6=2\times 3$ & $6=2\times 3$ \\
IX = 12--2B$_3$ & \text{none} & \text{none} & $8=4\text{ conics}$ & $8=4\text{ conics}$ \\
X = 12--B$_3$--C$_2$ & $3=1+\text{edge twice}$ & $3=1+\text{rec.\ of edge twice}$ & $6=2\times 3$ & $6=2\times 3$ \\
XI = 12--B$_5$ & $3=\text{edge 3 times}$ & $3=\text{rec.\ of edge 3 times}$ & $6=3\text{ conics}$ & $6=3\text{ conics}$ \\
XII = 12--U$_3$ & 3 & 3 & $6=2\times 3$ & $6=2\times 3$ \\
XIII = 12--B$_3$--2C$_2$ & 1 & 1 & $4=2\times 2$, nodal quadriquadric & $4=2\times 2$ quadriquadric \\
XIV = 12--B$_5$--C$_2$ & $1=\text{edge}$ & $1=\text{rec.\ of edge}$ & $4=3+\text{edge}$ & $4=3+\text{rec.\ of edge}$ \\
XV = 12--U$_5$ & 1 & 1 & $4=2\times 2$, nodal quadriquadric & $4=2\times 2$ cuspidal quadriquadric \\
XVI = 12--4C$_2$ & 3 & 3 & \text{none} & \text{none} \\
XVII = 12--2B$_3$--C$_2$ & \text{none} & \text{none} & $2=\text{conic}$ & $2=\text{conic}$ \\
XVIII = 12--B$_4$--2C$_2$ & $3=1+\text{edge twice}$ & $1+\text{rec.\ of edge twice}$ & \text{none} & \text{none} \\
XIX = 12--B$_5$--C$_3$ & $3=\text{axis 3 times}$ & $3=\text{rec.\ of axis 3 times}$ & \text{none} & \text{none} \\
XX = 12--U$_5$ & \text{none} & \text{none} & $2=\text{conic}$ & $2=\text{conic}$ \\
XXI = 12--3B$_3$ & \text{none} & \text{none} & \text{none} & \text{none} \\
\bottomrule
\end{tabular}
\end{center}

I pass now to the two cases of cubic scrolls.

\subsection*{Section XXII = S(1,1).\ Article No.\ 203.\ Equation $X^2W+Y^2Z=0$.}

\noindent\textbf{203.} As this is a scroll there is here no question of the 27 lines and 45 planes; there is a nodal line $X=0$, $Y=0$, ($b=1$) and a single directrix line, $Z=0$, $W=0$. The Hessian surface is $X^2Y^2=0$; the complete intersection with the cubic surface is made up of $X=0$, $Y=0$ (the nodal line) eight times, and of the lines $X=0$, $Z=0$, and $Y=0$, $W=0$, each twice. The reciprocal surface is $x^2z-y^2w=0$; viz.\ this is a like scroll, XXII = S$(1,1)$; $c'=0$, $b'=1$.

\subsection*{Section XXIII = S(1$\overline{1}$1).\ Article No.\ 204.\ Equation $X(XW+YZ)+Y^2=0$.}

\noindent\textbf{204.} This is also a scroll; there is a nodal line $X=0$, $Y=0$, and a single directrix line united therewith. The Hessian surface is $X^4=0$; the complete intersection with the cubic surface is $X=0$, $Y=0$ (the nodal line) twelve times. The reciprocal surface is $w(xw+yz)-z^2=0$; viz.\ this is a like scroll, XXIII = S$(1\overline{1}1)$; $c'=0$, $b'=1$.


\subsection*{Annex containing Additional Researches in regard to the case XX = 12--U$_{\mbox{\scriptsize 5}}$;\ equation $WX^2+XZ^2+Y^3=0$.}

\noindent\textbf{205.} Let the surface be touched by the line $(a,b,c,f,g,h)$, that is, the line the equations whereof are
\[
\begin{vmatrix}
0, & h, & -g, & a \\
-h, & 0, & f, & b \\
g, & -f, & 0, & c \\
-a, & -b, & -c, & 0
\end{vmatrix}
\!\hbox{\hspace{-0.5em}}\raisebox{0.3ex}{$\scriptstyle\frown$}\!(X,Y,Z,W)=0.
\]
Writing the equation in the form $cW\cdot cX^2+X(cZ)^2+c^2Y^3=0$, and substituting for $cW$, $cZ$ their values in terms of $X$, $Y$, we have $(-gX+fY)cZ^2+X(aX+bY)^2+c^2Y^3=0$, that is
\[
(3(a^2-cg),\ 2ab+cf,\ b^2,\ 3c^2\!\hbox{\hspace{-0.5em}}\raisebox{0.3ex}{$\scriptstyle\frown$}\!X,Y)^3=0,
\]
viz.\ the condition of contact is obtained by equating to zero the discriminant of the cubic function. We have thus
\begin{multline*}
27c^4(a^2-cg)^2 + 4c^3f^2(a^2-cg)^2 + c^2f^2(2ab+cf)^2 - b^4(2ab+cf)^2 \\
- 18b^2c^2(a^2-cg)(2ab+cf)=0,
\end{multline*}
viz.\ this is
\begin{multline*}
+27a^4c^4 - 4a^3b^2f + 30a^2b^2c^2f - 54a^2c^4g + 36ab^3cfg + 24abc^3f^2g \\
+ 4b^6f^2 - 18b^4c^2fg + 18b^2c^4f^2g + 27c^8g^2 + 4c^4f^3g^2=0,
\end{multline*}
which is the condition in order that the line $(a,b,c,f,g,h)$ may touch the surface $X^2W+XZ^2+Y^3=0$; and if we unite thereto the conditions that the line shall pass through a given point $(\alpha,\beta,\gamma,\delta)$, we have in effect the equation of the circumscribed cone, vertex $(\alpha,\beta,\gamma,\delta)$.

\setcounter{page}{452}

Writing $(f,g,h,a,b,c)$ in place of $(a,b,c,f,g,h)$, we obtain
\begin{multline*}
27f^4 - 4f^3g^2a + 30f^2g^2ha - 54f^2h^4b + 36fg^3h^2b + 24fgh^3a^2 \\
+ 4g^6a^2 - 18g^4h^2ab + 18g^2h^4a^2b + 27h^8b^2 + 4h^4a^3b^2=0
\end{multline*}
as the condition that the line $(a,b,c,f,g,h)$ shall touch the reciprocal surface $27(4x^2w+z^2)^2+64y^3w=0$; and if we consider $a$, $b$, $c$, $f$, $g$, $h$ as standing for
\[
\beta z-\gamma y,\ \gamma x-\alpha z,\ \alpha y-\beta x,\ \delta x-\alpha w,\ \delta y-\beta w,\ \delta z-\gamma w,
\]
values which satisfy the relation
\[
\begin{vmatrix}
0, & h, & -g, & a \\
-h, & 0, & f, & b \\
g, & -f, & 0, & c \\
-a, & -b, & -c, & 0
\end{vmatrix}
\!\hbox{\hspace{-0.5em}}\raisebox{0.3ex}{$\scriptstyle\frown$}\!(\alpha,\beta,\gamma,\delta)=0,
\]
then the equation in $(a,b,c,f,g,h)$ is that of the circumscribed cone, vertex $(\alpha,\beta,\gamma,\delta)$; the order being (as it should be) $a'=6$.

The cuspidal conic is $y=0$, $4x^2w+z^2=0$, and we at once obtain $a^2-4cg=0$ as the condition that the line $(a,b,c,f,g,h)$ shall pass through the cuspidal cone. Hence attributing to $(a,b,c,f,g,h)$ the foregoing values, we have $a^2-4cg=0$ for the equation of the cone, vertex $(\alpha,\beta,\gamma,\delta)$, which passes through the cuspidal conic; this is of course a quadric cone, $c'=2$.

\setcounter{page}{453}

I proceed to determine the intersections of the two cones. Representing by $\Theta=0$ the foregoing equation of the circumscribed cone, and putting for shortness
\[
X=27h^4(f^2-bh)-2g^2(2fg+ah),
\]
I find that we have identically
\[
\Theta - (f^2-bh)X + (g^2-4ah-8fgh)(a^2-4cg) - (32fg^2h+64agh^2)(af+bg+ch)=0:
\]
whence in virtue of the relation $af+bg+ch=0$, we see that the equations $\Theta=0$, $a^2-4cg=0$, are equivalent to $(f^2-bh)X=0$, $a^2-4cg=0$, or the twelve lines of intersection break up into the two systems
\begin{align*}
f^2-bh&=0, & a^2-4cg&=0, \quad \text{and} \\
X=27h^4(f^2-bh)-2g^2(2fg+ah)&=0, & a^2-4cg&=0.
\end{align*}

\setcounter{page}{454}

To determine the lines in question, observe that we have the relation between $(a,b,c,f,g,h)$ and $(\alpha,\beta,\gamma,\delta)$, and we can by the first three of these express $a$, $b$, $c$ linearly in terms of $f$, $g$, $h$; the equations $f^2-bh=0$, $a^2-4cg=0$, $27h^4(f^2-bh)-2g^2(2fg+ah)=0$ become thus homogeneous equations in $(f,g,h)$; the equations may in fact be written
\begin{align*}
8\delta^2(a^2-4cg) &= (\gamma^2+4\alpha\delta)g^2+\beta^2h^2-2\beta\gamma gh-4\beta\delta hf=0,\\
8(f^2-bh) &= 8f^2-\alpha h^2+\gamma hf=0,\\
8X &= 27h^4(8f^2-\alpha h^2+\gamma hf)+2g^2(\beta h^2-\gamma gh-2\delta fg)=0,
\end{align*}
viz.\ interpreting $(f,g,h)$ as coordinates in piano, the first equation represents a conic, the second a pair of lines, and the third a quartic.

\setcounter{page}{455}

We have identically
\[
\{2\beta\delta f-(\gamma^2+4\alpha\delta)g+\beta\gamma h\}^2-4\beta^2\delta(8f^2-\alpha h^2+\gamma hf) = (\gamma^2+4\alpha\delta)\{(\gamma^2+4\alpha\delta)g^2+\beta^2h^2-2\beta\gamma gh-4\beta\delta hf\};
\]
and it thus appears that the two conics touch at the points given by the equations $8f^2-\alpha h^2+\gamma hf=0$, $2\beta\delta f-(\gamma^2+4\alpha\delta)g+\beta\gamma h=0$: we have moreover
\[
-(\gamma^2+4\alpha\delta)(\beta h^2-\gamma gh-2\delta fg) = 4\beta\delta(8f^2-\alpha h^2+\gamma hf) + (-2\delta f-\gamma h)\{2\beta\delta f-(\gamma^2+4\alpha\delta)g+\beta\gamma h\},
\]
hence at the last-mentioned two points $-\beta h^2+\gamma gh+2\delta fg$ is $=0$; and the quartic $X=0$ thus passes through these two points.

The conic $(a^2-cg)=0$ and the quartic $X=0$ intersect besides (as is evident) in the point $g=0$, $h=0$ reckoning as two points, since it is a node of the quartic; and they must consequently intersect in four more points.

\setcounter{page}{456}

Recapitulating, the conic $a^2-4cg=0$ meets the sextic $(f^2-bh)X=0$ in the two points
\[
8f^2-\alpha h^2+\gamma hf=0, \quad 2\beta\delta f-(\gamma^2+4\alpha\delta)g+\beta\gamma h=0
\]
each three times, in the point $g=0$, $h=0$ twice, and in the four points
\begin{align*}
&160\delta g^4-27(\gamma^2+4\alpha\delta)g^2h^2+27\beta^2h^4=0,\\
&(\gamma^2+4\alpha\delta)g^2+\beta^2h^2-2\beta\gamma gh-4\beta\delta hf=0
\end{align*}
each once.

Or reverting to the proper significations of $(a,b,c,f,g,h)$, instead of points we have 2 lines each three times, a line twice, and 4 lines each once; the line $g=0$, $h=0$, that is, $g=0$, $h=0$, $a=0$, being, it will be observed, the line $\frac{x}{\alpha}=\frac{y}{\beta}=\frac{z}{\gamma}$ drawn from $(\alpha,\beta,\gamma,\delta)$ to the point $y=0$, $z=0$, $w=0$, which is the reciprocal of the uniplane $X=0$: the twelve lines are the $a'c'$ lines of intersection of the circumscribed cone $a'$ with the cuspidal cone $c'$, viz.\ $a'c'=[a'c']+3(\alpha'+\chi')$ where $[a'c']=4$ referring to the last-mentioned four lines; $\alpha'=2$ to the two lines; and $\chi'=2$ to the line $g=0$, $h=0$, $a=0$, which it thus appears must in the present case be reckoned twice.



%============================================================
% PAPER 413: A MEMOIR ON ABSTRACT GEOMETRY
%============================================================
\newpage
\section*{[413] A Memoir on Abstract Geometry.}
\markboth{A Memoir on Abstract Geometry.}{Cayley --- Abstract Geometry.}
\addcontentsline{toc}{section}{[413] A Memoir on Abstract Geometry}

\begin{flushright}
\small[From the \textit{Philosophical Transactions} of the Royal Society of London, vol.~CLX.\ (for the year 1870), pp.~51---63. Received October 14,---Read December 16, 1869.]
\end{flushright}

I SUBMIT to the Society the present exposition of some of the elementary principles of an Abstract $n$-dimensional Geometry. The science presents itself in two ways,---as a legitimate extension of the ordinary two- and three-dimensional geometries; and as a need in these geometries and in analysis generally. In fact whenever we are concerned with quantities connected together in any manner, and which are, or are considered as variable or determinable, then the nature of the relation between the quantities is frequently rendered more intelligible by regarding them (if only two or three in number) as the coordinates of a point in a plane or in space: for more than three quantities there is, from the greater complexity of the case, the greater need of such a representation; but this can only be obtained by means of the notion of a space of the proper dimensionality; and to use such representation, we require the geometry of such space.

An important instance in plane geometry has actually presented itself in the question of the determination of the number of the curves which satisfy given conditions: the conditions imply relations between the coefficients in the equation of the curve; and for the better understanding of these relations it was expedient to consider the coefficients as the coordinates of a point in a space of the proper dimensionality.

A fundamental notion in the general theory presents itself, slightly in plane geometry, but already very prominently in solid geometry; viz.\ we have here the difficulty as to the form of the equations of a curve in space, or (to speak more accurately) as to the expression by means of equations of the twofold relation between the coordinates of a point of such curve. The notion in question is that of a $k$-fold relation,---as distinguished from any system of equations (or onefold relations) serving for the expression of it, and as giving rise to the problem how to express such relation by means of a system of equations (or onefold relations). Applying to the case of solid geometry my conclusion in the general theory, it may be mentioned that I regard the twofold relation of a curve in space as being completely and precisely expressed by means of a system of equations $(P=0,\ Q=0,\ \ldots\ T=0)$, when no one of the functions $P$, $Q$, \ldots\ $T$ is a linear function, with constant or variable integral coefficients, of the others of them, and when every surface whatever which passes through the curve has its equation expressible in the form $U=AP+BQ\ldots+\lambda T$, with constant or variable integral coefficients $A$, $B$, \ldots\ $\lambda$.

It is hardly necessary to remark that all the functions and coefficients are taken to be rational functions of the coordinates, and that the word integral has reference to the coordinates.

\setcounter{page}{456}

\subsection*{Article Nos.\ 1 to 36.\ General Explanations; Relation, Locus, \&c.}

\noindent\textbf{1.} Any $m$ quantities may be represented by means of $m+1$ quantities as the ratios of $m$ of these to the remaining $(m+1)$th quantity, and thus in place of the absolute values of the $m$ quantities we may consider the ratios of $m+1$ quantities.

\noindent\textbf{2.} It is to be noticed that we are throughout concerned with the ratios of the $m+1$ quantities, not with the absolute values; this being understood, any mention of the ratios is in general unnecessary; thus I shall speak of a relation between the $m+1$ quantities, meaning thereby a relation as regards the ratios of the quantities; and so in other cases. It may also be noticed that in many instances a limiting or extreme case is sometimes included, sometimes not included, under a general expression; the general expression is intended to include whatever, having regard to the subject-matter and context, can be included under it.

\noindent\textbf{3.} \textit{Postulate.} We may conceive between the $m+1$ quantities a relation.

\noindent\textbf{4.} A relation is either regular, that is, it has a definite manifoldness, or, say, it is a $k$-fold relation; or else it is irregular, that is, composed of relations not all of the same manifoldness.

\noindent\textbf{5.} The ratios are determined (not in general uniquely) by means of a $m$-fold relation; and a relation cannot really be more than $m$-fold. But the notion of a more than $m$-fold relation has nevertheless to be considered. A relation may be, either in mere appearance or else according to a provisional conception thereof, more than $m$-fold, and be really $m$-fold or less than $m$-fold. Thus a relation expressed by $m+1$ or more equations is in general and prim\^{a} facie more than $m$-fold; but if the equations are not independent, and equivalent to $m$ or fewer equations, then the relation will be $m$-fold or less than $m$-fold. Or the given relation may depend on parameters, and so long as these are arbitrary be really more than $m$-fold; but the parameters may have to be, and be accordingly, so determined that the relation shall be $m$-fold or less than $m$-fold. A more than $m$-fold relation is said to be superdeterminate.

\noindent\textbf{6.} A system of any number of onefold relations, whether independent or dependent, and if more than $m$ of them, whether compatible or incompatible, is termed a `Plexus.

\noindent\textbf{7.} The number of onefold relations of a plexus may be $0$, $1$, $2$, \ldots\ or any number whatever; if it is $=m$, the plexus is said to be complete.

\noindent\textbf{8.} Any onefold relation whatever may be termed an equation; a system of onefold relations is a system of equations.

\noindent\textbf{9.} Any $k$ onefold relations form a $k$-fold relation; or say they are a $k$-fold relation.

\noindent\textbf{10.} A $k$-fold relation between $m+1$ quantities determines (not in general uniquely) the ratios of the quantities; a $k$-fold relation between the coordinates determines a locus of the dimensionality $m-k$, or say it is a $(m-k)$-dimensional locus; in particular a onefold relation determines a $(m-1)$-dimensional locus, which may be termed an `Omal.

\noindent\textbf{11.} A $k$-fold relation is implied in a $l$-fold relation when every set of values of the coordinates which satisfies the $l$-fold relation satisfies also the $k$-fold relation. In particular a onefold relation is implied in a $k$-fold relation when every set of values which satisfies the $k$-fold relation satisfies also the onefold relation.

\noindent\textbf{12.} The foregoing definition of implication is applicable to the case of a $k$-fold relation and a $l$-fold relation where $k>l$, or even where $k=l$; but observe that a $k$-fold relation is not in general implied in a different $k$-fold relation.

\noindent\textbf{13.} If a $k$-fold relation is implied in each of two or more $l$-fold relations ($k$ being $>l$ or $=l$), then the $k$-fold relation is said to be implied in the aggregate of the $l$-fold relations. And more generally, if a $k$-fold relation is implied in each of two or more relations not all of them of the same manifoldness, then it is implied in the aggregate of these relations.

\setcounter{page}{457}

\noindent\textbf{14.} Any two or more relations may be composed together, and they are then factors of a single composite relation; viz.\ the composite relation is a relation satisfied if, and not satisfied unless, some one of the component relations be satisfied.

\noindent\textbf{15.} The foregoing notion of composition is, it will be noticed, altogether different from that which would at first suggest itself. The definition is defective as not explaining the composition of a relation any number of times with itself, or elevation thereof to power; which however must be admitted as part of the notion of composition.

\noindent\textbf{16.} A $k$-fold relation which is not satisfied by any other $k$-fold relation, and which is not a power, is a prime relation. A relation which is not prime is composite, viz.\ it is a relation composed of prime factors each taken once or any other number of times; in particular, it may be the power of a single prime factor. Any prime factor is single or multiple according as it occurs once or a greater number of times.

\noindent\textbf{17.} A relation which is either prime, or else composed of prime factors each of the same manifoldness, is a regular relation; a $k$-fold relation is ex vi termini regular. An irregular relation is a composite relation the prime factors whereof are not all of the same manifoldness.

\noindent\textbf{18.} A prime $k$-fold relation cannot be implied in any prime $k$-fold relation different from itself. But a prime $k$-fold relation may be implied in a prime more-than-$k$-fold relation,---or in a composite relation, regular or irregular, each factor whereof is more than $k$-fold; and so also a composite relation, regular or irregular, each factor whereof is at most $k$-fold, may be implied in a composite relation, regular or irregular, each factor whereof is more than $k$-fold. In a somewhat different sense, each factor of a composite relation implies the composite relation.

\noindent\textbf{19.} A composite relation is satisfied if any particular one of the component relations is satisfied; but in order to exclude this case we may speak of a composite relation as being satisfied distributively; viz.\ this will be the case if, in order to the satisfaction of the composite relation, it is necessary to consider all the factors thereof, or, what is the same thing, when the reduced relation obtained by the omission of any one factor whatever is not always satisfied. And when the composite relation is satisfied distributively, the several factors thereof are satisfied alternatively; viz.\ there is no one which is throughout unsatisfied.

\noindent\textbf{20.} A composite onefold relation is never distributively implied in a prime $k$-fold relation---that is, a prime $k$-fold relation implies only a prime onefold relation; or at least only implies a composite onefold relation improperly, in the sense that it implies a certain prime factor of such composite onefold relation. Conversely, every $k$-fold relation which implies distributively a composite onefold relation is composite.

\noindent\textbf{21.} Any two or more relations may be aggregated together, and they are then constituents of a single aggregate relation; viz.\ the aggregate relation is a relation satisfied if, and only if, all the constituent relations are satisfied. The aggregation of relations corresponds to the intersection of loci.

\noindent\textbf{22.} There is no meaning in aggregating a relation with itself; such aggregation only occurs accidentally when two relations aggregated together become one and the same relation; and the aggregate of a relation with itself is nothing else than the original relation.

\noindent\textbf{23.} A onefold relation is not an aggregate, but is its own sole constituent; a more than onefold relation may always be considered as an aggregate of two or more constituent relations. The constituent relations determine, they in fact constitute, the aggregate relation; but the aggregate relation does not in any wise determine the constituent relations. Any relation implied in a given relation may be considered as a constituent of such given relation.

\setcounter{page}{458}

\noindent\textbf{24.} The aggregate of a $k$-fold and a $l$-fold relation is in general and at most a $(k+l)$-fold relation; when it is a $(k+l)$-fold relation, the constituent relations are independent, but otherwise, viz.\ if the aggregate relation is, or has for factor, a less than $(k+l)$-fold equation, the constituent relations are dependent or interconnected.

\noindent\textbf{25.} Passing from relations to loci, we may say that the composition of relations corresponds to the congregation of loci, and the aggregation of relations to the intersection of loci.

\noindent\textbf{26.} For, first, the locus (if any) corresponding to a given composite relation is the congregate of the loci corresponding to the several prime factors of the given relation, the locus corresponding to a single factor being taken once, and the locus corresponding to a multiple factor being taken a number of times equal to the multiplicity of the factor.

\noindent\textbf{27.} And, secondly, the locus (if any) corresponding to a given aggregate relation is the locus common to and contained in each of the loci corresponding to the several constituent relations respectively; or, what is the same thing, it is the intersection of these several loci.

\noindent\textbf{28.} It may be remarked that a $k$-fold locus and a $l$-fold locus where $k+l>m$ (or where the aggregate relation is more than $m$-fold) have not in general any common locus.

\noindent\textbf{29.} Any onefold relation implied in a given $k$-fold relation is said to be in involution with the $k$-fold relation, and so in a system of onefold relations, if any relation be implied in the other relations, or, what is the same thing, in the relation aggregated of the other relations, then the system is said to be in involution; a system not in involution is said to be asyzygetic.

\noindent\textbf{30.} Consider a given $k$-fold relation, and, in conjunction therewith, a system of any number of onefold relations each implied in the given $k$-fold relation. We may omit from the system any relation implied in the remaining relations, and so successively until we arrive at an asyzygetic system. Consider now any other onefold relation implied in the given $k$-fold relation; this is either implied in the system of onefold relations, and it is then to be rejected, or if it is not implied in the system, it is to be added on to and made part of the system. It may happen that, in the system thus obtained, some one relation of the original system is implied in the remaining relations of the new system; but if this is so the implied relation is to be rejected; the new system will in this case contain only as many relations as the original system, and in any case the new system will be asyzygetic. Treating in the same manner every other onefold relation implied in the given $k$-fold relation, we ultimately arrive at an asyzygetic system of onefold relations, such that every onefold relation implied in the given $k$-fold relation is implied in the asyzygetic system. The number of onefold relations will be at least equal to $k$ (for if this were not so we should have the given $k$-fold relation as an aggregate of less than $k$ onefold relations); but it may be greater than $k$, and it does not appear that there is any [assignable] superior limit to the number of onefold relations of the asyzygetic system.

\setcounter{page}{459}

\noindent\textbf{31.} The system of onefold relations is a precise equivalent of the given $k$-fold relation. Every set of values of the coordinates which satisfies the given $k$-fold relation satisfies the system of onefold relations; and reciprocally every set of values which satisfies the system of onefold relations satisfies the given $k$-fold relation. But if we omit any one or more of the onefold relations, then the reduced system so obtained is not a precise equivalent of the given $k$-fold relation; viz.\ there exist sets of values satisfying the reduced system, but not satisfying the given $k$-fold relation.

\noindent\textbf{32.} In fact consider a $k$-fold relation the aggregate of less than all of the onefold relations of the asyzygetic system, and in connexion therewith an omitted onefold relation; this omitted relation is not implied in the aggregate, and it constitutes with the aggregate not a $(k+1)$-fold, but only a $k$-fold relation. This happens as follows, viz.\ the omitted relation is a factor of a composite onefold relation distributively implied in the aggregate; hence the aggregate is composite, and it implies distributively a composite onefold relation composed of the omitted relation and of an associated onefold relation; that is, the aggregate will be satisfied by values which satisfy the omitted relation, and also by values which (not satisfying the omitted relation) satisfy the associated relation just referred to.

\noindent\textbf{33.} Selecting at pleasure any $k$ of the onefold relations of the asyzygetic system, being such that the aggregate of the $k$ relations is a $k$-fold relation, we have a composite $k$-fold relation wherein each of the remaining onefold relations is alternatively implied; viz.\ each remaining onefold relation is a factor of a composite onefold relation implied distributively in the composite $k$-fold relation. Hence considering the $k+1$ onefold relations, viz.\ any $k+1$ relations of the asyzygetic system, each one of these is implied alternatively in the aggregate of the other $k$ relations.

\noindent\textbf{34.} The asyzygetic system of onefold relations thus possesses the property that no reduced system thereof (that is, system obtained by the omission of one or more relations from the asyzygetic system) is a precise equivalent of the given $k$-fold relation; and moreover any relation alternatively implied in the asyzygetic system is one of the relations of the system. It may be added that, besides the relations of the system, there is not any onefold relation alternatively implied in the asyzygetic system.

\noindent\textbf{35.} The foregoing theory has been stated without any limitation as to the value of $k$, and it has I think a meaning even when $k$ is $>m$; but the ordinary case is $k\pitchfork m$. Considering the theory as applying to this case, I remark that the last proposition, viz.\ that no reduced system is a precise equivalent of the given $k$-fold relation, is generally true only on the assumption of the existence or quasi-existence of sets of values satisfying a more than $m$-fold relation. For let $k$ be $>m$, and, on the contrary, assume, as we usually do, that it is not in general possible to satisfy a more than $m$-fold relation between the coordinates; the number of relations in the system may be $>m+1$; and if this is so, then selecting any $m+1$ relations of the system, it may very well happen that the given $k$-fold relation is not satisfied by any sets of values other than those which satisfy the $m+1$ relations,---that is, that the $m+1$ relations are a precise equivalent of the given $k$-fold relation. But even in this case the consideration of the entire system of the onefold relations is not the less advantageous; and I say in general that the given $k$-fold relation has for its precise and complete equivalent the asyzygetic system of onefold relations.

\setcounter{page}{460}

\noindent\textbf{36.} (In illustration of the foregoing Nos.\ 29 to 35, I remark that, for the functions or equations $P=0$, $Q=0$, $R=0$, \&c., if we have identically $AP+BQ+CR+\ldots=0$, where the factors $A$, $B$, $C$, \ldots\ are integral functions of the coordinates, and where some one of these factors, say, $A$, is a constant (or if we please $=1$), then the system of functions or equations is in involution; or, to speak more accurately, the function or equation $P=0$ is in involution with the remaining functions or equations $Q=0$, $R=0$, \ldots\ But when the factors $A$, $B$, $C$, \ldots\ are no one of them constant, then we have a convolution. If $P=0$ is in involution with the remaining equations $Q=0$, $R=0$, \ldots, then $P=0$ is implied in these equations, and the relations $(Q=0,\ P=0,\ \ldots)$ and $(P=0,\ Q=0,\ R=0,\ \ldots)$ are equivalent to each other. But in the case of a convolution where $AP+BQ+CR+\ldots=0$, then the relation the equations $Q=0$, $R=0$, \ldots\ imply $AP=0$, that is, $A=0$ or else $P=0$; or, what is the same thing, the relation $(Q=0,\ R=0,\ \ldots)$ is a relation composed of the two relations $(A=0,\ Q=0,\ R=0,\ \ldots)$ and $(P=0,\ Q=0,\ R=0,\ \ldots)$. In the $k$-fold relation expressed by the more than $k$ equations $(P=0,\ Q=0,\ R=0,\ \ldots)$, selecting any $k$ of these equations which are not in convolution, and uniting thereto any one of the remaining equations, we have a convolution of $k+1$ equations; and when a $k$-fold relation is precisely expressed by means of a system of $k$ or more equations $(P=0,\ Q=0,\ \ldots)$, then every equation $U=0$ implied in the given relation, or, what is the same thing, the aggregate relation $(P=0,\ Q=0,\ \ldots,\ U=0)$, is a system in involution.)

\subsection*{Article Nos.\ 37 to 42.\ Omal Relation; Order.}

\setcounter{page}{461}

\noindent\textbf{37.} A $k$-fold relation may be linear or omal. If $k=m$, the corresponding locus is a point; if $k<m$ the locus is a $k$-fold, or $(m-k)$-dimensional omaloid; the expression omaloid used absolutely denotes the onefold or $(m-1)$-dimensional omaloid; the point may be considered as a $m$-fold omaloid.

\noindent\textbf{38.} A $m$-fold relation which is not linear or omal is of necessity composite, composed of a certain number $M$ of $m$-fold linear or omal relations; viz.\ the $m$-fold locus corresponding to the $m$-fold relation is a point-system of $M$ points, each of which may be considered as given by a separate $m$-fold linear or omal relation; each which relation is a factor of the original $m$-fold relation. The given $m$-fold relation, and the point-system corresponding thereto, are respectively said to be of the order $M$.

\noindent\textbf{39.} The order of a point-system of $M$ points is thus $=M$, but it is of course to be borne in mind that the points may be single or multiple points; and that if the system consists of a point taken $\alpha$ times, another point taken $\beta$ times, \&c., then the number of points and therefore the order $M$ of the system is considered to be $=\alpha+\beta+\ldots$.

\noindent\textbf{40.} If to a given $k$-fold relation ($k<m$) we unite an absolutely arbitrary $(m-k)$-fold linear relation, so as to obtain for the aggregate a $m$-fold relation, then the order $M$ of this $m$-fold relation (or, what is the same thing, the number $M$ of points in the corresponding point-system) is said to be the order of the given $k$-fold relation. The notion of order does not apply to a more than $m$-fold relation.

\noindent\textbf{41.} The foregoing definition of order may be more compendiously expressed as follows: viz.\ Given between the $m+1$ coordinates a relation which is at most $m$-fold; then if it is not $m$-fold, join to it an arbitrary linear relation so as to render it $m$-fold; we have a $m$-fold relation giving a point-system; and the order of the given relation is equal to the number of points of the point-system.

\noindent\textbf{42.} The relation aggregated of two or more given relations, when the notion of order applies to the aggregate relation, that is, when it is not more than $m$-fold, is of an order equal to the product of the orders of the constituent relations; or, say, the orders of the given relations being $\mu$, $\mu'$, \ldots, the order of the aggregate relation is $=\mu\mu'\ldots$.

\subsection*{Article Nos.\ 43 and 44.\ Parametric Relations.}

\setcounter{page}{462}

\noindent\textbf{43.} We have considered so far relations which involve only the coordinates $(x,y,\ldots)$; the coefficients are purely numerical, or, if literal, they are absolute constants, which either do or do not satisfy certain conditions; if they do not, the relation assumed in the first instance to be $k$-fold is really $k$-fold, or, as we may express it, the relation is really as well as formally $k$-fold; if they do satisfy certain relations in virtue whereof the formally $k$-fold relation is really less than $k$-fold, say, it is $(k-l)$-fold, then the relation is in fact to be considered ab initio as a $(k-l)$-fold relation: there is no question of a relation being in general $k$-fold and becoming less than $k$-fold, or suffering any other modification in its form; and the notion of a more than $m$-fold relation is in the preceding theory meaningless.

\noindent\textbf{44.} But a relation between the coordinates $(x,y,\ldots)$ may involve parameters, and so long as these remain arbitrary it may be really as well as formally $k$-fold; but when the parameters satisfy certain conditions, it may become $(k-l)$-fold, or may suffer some other modification in its form. And we have to consider the theory of a relation between the coordinates $(x,y,\ldots)$, involving besides parameters which may satisfy certain conditions, or, say simply, a relation involving variable parameters. If the number of the parameters be $m'$, then these parameters may be regarded as the ratios of $m'$ quantities to a remaining $m'+1$th quantity, and the relation may be considered as involving homogeneously the $m'+1$ parameters $(x',y',\ldots)$. And these may, if we please, be regarded as coordinates of a point in their own $m'$-dimensional space, or we have to consider relations between the $m+1$ coordinates $(x,y,\ldots)$ and the $m'+1$ (parameters or) coordinates $(x',y',\ldots)$.

\subsection*{Article Nos.\ 45 to 55.\ Quantics, Notation, \&c.}

\setcounter{page}{463}

\noindent\textbf{45.} A homogeneous function of the coordinates $(x,y,\ldots)$ is represented by a notation such as $({*}\!\hbox{\hspace{-0.5em}}\raisebox{0.3ex}{$\scriptstyle\frown$}\!x,y,\ldots)^r$ (where $(*)$ indicates the coefficients and $(r)$ the degree), and it is said to be a quantic; and in reference to the quantic the quantities or coordinates $(x,y,\ldots)$ are also termed facients. More generally a quantic involving two or more sets of coordinates, or facients, is represented by the similar notation $({*}\!\hbox{\hspace{-0.5em}}\raisebox{0.3ex}{$\scriptstyle\frown$}\!x,y,\ldots)^r(x',y',\ldots)^{r'}\ldots$.

\noindent\textbf{46.} The quantic is unipartite, bipartite, tripartite, \&c., according as the number of sets is one, two, three, \&c.; and with respect to any set of coordinates, it is binary, ternary, quaternary, \ldots\ $(m+1)$ary, according as the number of the coordinates is two, three, four, or $m+1$; and it is linear, quadric, cubic, quartic, \ldots, according as the degree in regard to the coordinates in question is 1, 2, 3, 4, \ldots.

\noindent\textbf{47.} A quantic involving two or more sets of coordinates, and linear in regard to each of them, is said to be multipartite; or, in particular, when there are only two sets, it is said to be lineo-linear; we may even extend the epithet lineo-linear to the case of any number of sets.

\noindent\textbf{48.} Instead of the general notation we may write $(a,\ldots\!\hbox{\hspace{-0.5em}}\raisebox{0.3ex}{$\scriptstyle\frown$}\!x,y,\ldots)^r(x',y',\ldots)^{r'}\ldots$, where the coefficients are now indicated by $(a,\ldots)$, and the degrees are $r$, $r'$, \ldots.

\noindent\textbf{49.} In the cases where the particular values of the coefficients have to be attended to, we write down the entire series of coefficients, or at least refer thereto by the notation $(a,\ldots)$; and it is to be understood that the coefficients expressed or referred to are numerical or else they are independent arbitrary quantities.

\noindent\textbf{50.} In the case of a binary quantic, the series of coefficients is $a$, $b$, \ldots\ $h$, $k$, $l$, or as we may write them $(a,b,\ldots,l)$, and the number of the coefficients is $=r+1$, if $r$ is the degree. In the case of a ternary quantic, the series of coefficients may be written $(a,\ldots\pitchfork\!\hbox{\hspace{-0.5em}}\raisebox{0.3ex}{$\scriptstyle\frown$}\!x,y,z)^r$.

\noindent\textbf{51.} For shortness, the quantic may be represented by a single letter $U$; and when this is so, the equation $U=0$ is used to denote the onefold relation between the coordinates expressed by the quantic being $=0$.

\noindent\textbf{52.} A onefold relation between the coordinates is expressible by means of an equation of the form $({*}\!\hbox{\hspace{-0.5em}}\raisebox{0.3ex}{$\scriptstyle\frown$}\!x,y,\ldots)^r=0$.

\noindent\textbf{53.} The expression ``an equation'' used without explanation may be taken to mean an equation of the form in question, viz.\ the equation obtained by putting a quantic equal to zero; the quantic is said to be the nilfactum of the equation.

\noindent\textbf{54.} It is frequently convenient to denote the quantic or nilfactum by a single letter, and to use a locution such as ``the equation $U=({*}\!\hbox{\hspace{-0.5em}}\raisebox{0.3ex}{$\scriptstyle\frown$}\!x,y,\ldots)^r=0$,'' which really means that the single letter $U$ stands for the quantic $({*}\!\hbox{\hspace{-0.5em}}\raisebox{0.3ex}{$\scriptstyle\frown$}\!x,y,\ldots)^r$, so that we are afterwards at liberty to write $U=0$ as an abbreviated expression for $({*}\!\hbox{\hspace{-0.5em}}\raisebox{0.3ex}{$\scriptstyle\frown$}\!x,y,\ldots)^r=0$.

\noindent\textbf{55.} A $k$-fold relation between the coordinates is (as has been shown) equivalent to a system of $k$ or more onefold relations; each of these is expressible by an equation $U=0$, and the $k$-fold relation is thus expressible by a system of $k$ or more such equations. Representing by $((U))$ the system of functions which are the nilfacta of these equations respectively, the $k$-fold relations may be represented thus, $((U))=0$; or more completely, the relation being $k$-fold, and the number of equations being $=s$, by the notation $\{((U)_s)\}^{(k\text{-fold})}=0$.

\subsection*{Article Nos.\ 56 to 62.\ Resultant, Discriminant, \&c.}

\setcounter{page}{466}

\noindent\textbf{56.} In the case $k>m$, a given $k$-fold relation between the $m+1$ coordinates $(x,y,\ldots)$ and the parameters $(x',y',\ldots)$ leads to a $(k-m)$-fold relation between the parameters. This is termed the resultant relation of the given $k$-fold relation, or when the additional specification is necessary, the resultant relation obtained by elimination of the coordinates $(x,y,\ldots)$.

\noindent\textbf{57.} Consider a $k$-fold relation between the $m+1$ coordinates $(x,y,\ldots)$ and the $m'+1$ coordinates $(x',y',\ldots)$. If $k\pitchfork m$, then, considering the $(x,y,\ldots)$ as coordinates and the $(x',y',\ldots)$ as parameters, we have corresponding to the given relation a $k$-fold locus in the $m$-space; and so if $k\pitchfork m'$, then, considering the $(x',y',\ldots)$ as coordinates, but the $(x,y,\ldots)$ as parameters, we have corresponding to the given relation a $k$-fold locus in the $m'$-space.

\noindent\textbf{58.} If $k>m$, but if the $(k-m)$-fold resultant relation is satisfied, then the given $k$-fold relation becomes a $m$-fold linear relation between the coordinates $(x,y,\ldots)$, and is consequently satisfied by a single set of values of the coordinates. Hence, considering the given $k$-fold relation as implying the $(k-m)$-fold resultant relation, the $k$-fold relation will represent a single point in the $m$-space, say, the common point.

\noindent\textbf{59.} A $m$-fold relation, or the locus, or point-system thereby represented, may have a double or nodal point, viz.\ two of the points of the point-system may be coincident. More generally a $k$-fold relation ($k\pitchfork m$), or the locus thereby represented, may have a double or nodal point; for let the relation if less than $m$-fold be made $m$-fold by adjoining to it a linear $(m-k)$-fold relation satisfied by the coordinates of the point in question but otherwise arbitrary, then, if the point in question be a double or nodal point of the $m$-fold relation, or of the point-system thereby represented, the point is said to be a double or nodal point of the original $k$-fold relation, or of the locus thereby represented.

\noindent\textbf{60.} A given $k$-fold relation ($k\pitchfork m$) between the $m+1$ coordinates, or the locus thereby represented, has not in general a nodal point. But if the relation involve the $m'+1$ parameters $(x',y',\ldots)$, then, if a certain onefold relation be satisfied between the parameters, there will be a nodal point. The onefold relation between the parameters is the discriminant relation of the given $k$-fold relation.

\noindent\textbf{61.} In the case in question, $k\pitchfork m$, the discriminant relation is the resultant relation of a $(m+1)$-fold relation which is the aggregate of the given $k$-fold relation with a certain relation called the Jacobian relation, or when the distinction is required, the Jacobian relation in regard to the $(x,y,\ldots)$.

\noindent\textbf{62.} Consider a $k$-fold relation ($k\pitchfork m$, $k\pitchfork m'$) between the $m+1$ coordinates $(x,y,\ldots)$ and the $m'+1$ coordinates $(x',y',\ldots)$. It has been seen that to a given set of values of the $(x',y',\ldots)$ or, say, to a given point in the $m'$-space, there corresponds a $k$-fold locus in the $m$-space, and that to a given set of values of the $(x,y,\ldots)$, or to a given point in the $m$-space, there corresponds a $k$-fold locus in the $m'$-space. The $k$-fold locus in the $m'$-space may have a nodal point; this will be the case if there is satisfied between the $(x,y,\ldots)$ a certain onefold relation, the discriminant relation of the given $k$-fold relation in regard to the $(x',y',\ldots)$. This onefold relation represents in the $m$-space a onefold locus, the envelope of the $k$-fold loci in the $m$-space corresponding to the several points of the $m'$-space. The property of the envelope is that to each point thereof there corresponds in the $m'$-space a $k$-fold locus having a nodal point.

\subsection*{Article Nos.\ 63--69.\ Consecutive Points; Tangent Omals.}

\setcounter{page}{468}

\noindent\textbf{63.} As the notions of proximity and remoteness have been thus far excluded, it is necessary to introduce them for the purpose of defining a tangent omal.

\noindent\textbf{64.} If the coordinates of the given point are $(x,y,\ldots)$, those of the consecutive point may be assumed to be $(x+\delta x,\ y+\delta y,\ \ldots)$, where $\delta x$, $\delta y$, \ldots\ are indefinitely small in regard to $(x,y,\ldots)$.

\noindent\textbf{65.} It may be remarked that, taking the coordinates to be $(x+X,\ y+Y,\ldots)$, there is no obligation to have $(X,Y,\ldots)$ indefinitely small; in fact whatever the magnitudes of these quantities are, if only $X:Y:\ldots=x:y:\ldots$, then the point $(x+X,\ y+Y,\ldots)$ will be the very same with the original point, and it is therefore clear that a consecutive point may be represented in the same manner with magnitudes, however large, of $X$, $Y$, \ldots\ But we may assume them indefinitely small, that is, the ratios $x+\delta x:y+\delta y:\ldots$, where $\delta x$, $\delta y$, \ldots\ are indefinitely small in regard to $(x,y,\ldots)$, will represent any set of ratios indefinitely near to the ratios $(x:y:\ldots)$. The foregoing quantities $(\delta x,\ \delta y,\ \ldots)$ are termed the increments.

\noindent\textbf{66.} Consider a $k$-fold relation between the $m+1$ coordinates $(x,y,\ldots)$, $k\pitchfork m$; the increments $(\delta x,\ \delta y,\ \ldots)$ are connected by a linear $k$-fold relation. The linear $k$-fold relation is satisfied if we assume the increments proportional to the coordinates---this is, in fact, assuming that the point remains unaltered. We may write $(\delta x,\ \delta y,\ \ldots)=(x,\ y,\ \ldots)$, since in such an equation only the ratios are attended to. But it may be preferable to write $(\delta x,\ \delta y,\ \ldots)=\lambda(x,\ y,\ \ldots)$. In particular if $k=m$, then the increments are connected by a linear $m$-fold relation; that is, the ratio of the increments is uniquely determined; and as the relation is satisfied by taking the increments proportional to the coordinates, it is clear that the values which the linear $m$-fold relation gives for the increments are in fact proportional to the coordinates: viz.\ there is not in this case any consecutive point.

\noindent\textbf{67.} Considering the $k$-fold relation as belonging to a $k$-fold locus in the $m$-space, so that $(x,y,\ldots)$ are the coordinates of a point on this locus, then if in the linear $k$-fold relation between the increments these increments are replaced by the coordinates $(x,y,\ldots)$ of a point in the $m$-space, then considering the original coordinates $(x,y,\ldots)$ as parameters, the locus of the point $(x,y,\ldots)$ is a $k$-fold omal locus: it is to be observed that, by what precedes, the linear $k$-fold relation is satisfied by writing therein the values $x:y:\ldots=x:y:\ldots$, that is, the $k$-fold omal locus passes through the original point $(x,y,\ldots)$; the $k$-fold omal locus is said to be the tangent-omal of the original $k$-fold locus at the point $(x,y,\ldots)$, which point is said to be the point of contact.

\setcounter{page}{469}

\noindent\textbf{68.} If in the original $k$-fold locus we replace $(x,y,\ldots)$ by $(x,y,\ldots)$, and combine therewith the $k$-fold linear relation, we have between the coordinates $(x,y,\ldots)$ a $2k$-fold relation (containing as parameters the coordinates $(x,y,\ldots)$); these parameters satisfy the original $k$-fold relation, and in virtue thereof the $2k$-fold relation breaks up into the two $k$-fold relations which correspond to the two factors of the composite onefold relation distributively implied in the aggregate of the $k$-fold locus and the $k$-fold tangent-omal; that is, it breaks up into the $k$-fold relation which gives the ineunt-point and the $k$-fold relation which gives the tangent-omal; in other words, the two $k$-fold loci intersect in the $k$-fold locus which is the $k$-fold locus of contact.

\noindent\textbf{69.} We thus arrive at the notion of the double generation of a $k$-fold locus, viz.\ such locus is the locus of the points, or, say, of the ineunt-points thereof; and it is also the envelope of the tangent-omals thereof. We have thus a theory of duality; I do not at present attempt to develope the theory, but it is necessary to refer to it, in order to remark that this theory is essential to the systematic development of a $m$-dimensional geometry; the original classification of loci as onefold, twofold, \ldots\ $(m-1)$-fold is incomplete, and must be supplemented with the loci reciprocally connected with these loci respectively. And moreover the theory of the singularities of a locus can only be systematically established by means of the same theory of duality; the singularities in regard to the ineunt-point must be treated of in connexion with the singularities in regard to the tangent-omal. These theories (that is, the classification of loci, and the establishment and discussion of the singularities of each kind of locus), vast as their extent is, should in the logical order precede that for which other reasons it may be expedient next to consider, the theory of Transformation, as depending on relations involving simultaneously the $m+1$ coordinates $(x,y,\ldots)$ and the $m'+1$ coordinates $(x',y',\ldots)$.



%============================================================
% PAPER 414: ON POLYZOMAL CURVES
%============================================================
\newpage
\section*{[414] On Polyzomal Curves, otherwise the Curves $\sqrt{lU}+\sqrt{mV}+\&c.=0$.}
\markboth{On Polyzomal Curves.}{Cayley --- Polyzomal Curves.}
\addcontentsline{toc}{section}{[414] On Polyzomal Curves}

\begin{flushright}
\small[From the \textit{Transactions} of the Royal Society of Edinburgh, vol.~xxv.\ (1868), pp.~1---110. Read 16th December 1867.]
\end{flushright}

If $U$, $V$, \&c., are rational and integral functions $({*}\!\hbox{\hspace{-0.5em}}\raisebox{0.3ex}{$\scriptstyle\frown$}\!x,y,z)^r$, all of the same degree $r$, in regard to the coordinates $(x,y,z)$, then $\sqrt{lU}+\sqrt{mV}+\&c.$ is a polyzome, and the curve $\sqrt{lU}+\sqrt{mV}+\&c.=0$ a polyzomal curve. Each of the curves $\sqrt{lU}=0$, $\sqrt{mV}=0$, \&c.\ (or say the curves $U=0$, $V=0$, \&c.) is, on account of its relation of circumscription to the curve $\sqrt{lU}+\sqrt{mV}+\&c.=0$, considered as a girdle thereto (\textit{z\={o}ma}), and we have thence the term ``zome'' and the derived expressions ``polyzome,'' ``zomal,'' \&c.

If the number of the zomes $\sqrt{lU}$, $\sqrt{mV}$, \&c.\ be $=\nu$, then we have a $\nu$-zome, and corresponding thereto a $\nu$-zomal curve; the curves $U=0$, $V=0$, \&c., are the zomal curves or zomals thereof. The cases $\nu=1$, $\nu=2$, are not, for their own sake, worthy of consideration; it is in general assumed that $\nu$ is $=3$ at least. It is sometimes convenient to write the general equation in the form $\sqrt{U}+\&c.=0$, where $l$, \&c.\ are constants.

The Memoir contains researches in regard to the general $\nu$-zomal curve; the branches thereof, the order of the curve, its singularities, class, \&c.; also in regard to the $\nu$-zomal curve $\sqrt{l(\Theta+L\Phi)}+\sqrt{m(\Theta+M\Phi)}+\&c.=0$, where the zomal curves $\Theta+L\Phi=0$, \&c., all pass through the points of intersection of the same two curves $\Theta=0$, $\Phi=0$ of the orders $r$ and $r-s$ respectively; included herein we have the theory of the depression of order as arising from the ideal factor or factors of a branch or branches.

A general theorem is given of ``the decomposition of a tetrazomal curve,'' viz.\ if the equation of the curve be $\sqrt{lU}+\sqrt{mV}+\sqrt{nW}+\sqrt{pT}=0$; then if $U$, $V$, $W$, $T$ are in involution, that is, connected by an identical equation $aU+bV+cW+dT=0$, and if $l$, $m$, $n$, $p$, satisfy the condition $\frac{a}{l}+\frac{b}{m}+\frac{c}{n}+\frac{d}{p}=0$, the tetrazomal curve breaks up into two trizomal curves, each expressible by means of any three of the four functions $U$, $V$, $W$, $T$; for example, in the form $\sqrt{l'U}+\sqrt{m'V}+\sqrt{p'T}=0$.

If, in this theorem, we take $p=0$, then the original curve is the trizomal $\sqrt{lU}+\sqrt{mV}+\sqrt{nW}=0$, $T$ is any function $=-\frac{1}{d}(aU+bV+cW)$, where, considering $l$, $m$, $n$ as given, $a$, $b$, $c$ are quantities subject only to the condition $\frac{a}{l}+\frac{b}{m}+\frac{c}{n}=0$, and we have the theorem of ``the variable zomal of a trizomal curve,'' viz.\ the equation of the trizomal $\sqrt{lU}+\sqrt{mV}+\sqrt{nW}=0$, may be expressed by means of any two of the three functions $U$, $V$, $W$, and of a function $T$ determined as above, for example in the form $\sqrt{l'U}+\sqrt{m'V}+\sqrt{n'T}=0$; whence also it may be expressed in terms of three new functions $T$, determined as above.

This theorem, which occupies a prominent position in the whole theory, was suggested to me by Mr Casey's theorem, presently referred to, for the construction of a bicircular quartic as the envelope of a variable circle.

\setcounter{page}{470}

In the $\nu$-zomal curve $\sqrt{l(\Theta+L\Phi)}+\&c.=0$, if $\Theta=0$ be a conic, $\Phi=0$ a line, the zomals $\Theta+L\Phi=0$, \&c.\ are conics passing through the same two points $\Theta=0$, $\Phi=0$, and there is no real loss of generality in taking these to be the circular points at infinity---that is, in taking the conics to be circles. Doing this, and using a special notation $\mathfrak{A}^\circ=0$ for the equation of a circle having its centre at a given point $A$, and similarly $\mathfrak{A}=0$ for the equation of an evanescent circle, or say of the point $A$, we have the $\nu$-zomal curve $\sqrt{l}\mathfrak{A}^\circ+\&c.=0$, and the more special form $\sqrt{l}\mathfrak{A}+\&c.=0$.

As regards the last-mentioned curve, $\sqrt{l}\mathfrak{A}+\&c.=0$, the point $A$ to which the equation $\mathfrak{A}=0$ belongs, is a focus of the curve, viz.\ in the case $\nu=3$, it is an ordinary focus, and in the case $\nu>3$, it is a special kind of focus, which, if the term were required, might be called a foco-focus; the Memoir contains an explanation of the general theory of the foci of plane curves.

For $\nu=3$, the equation $\sqrt{l}\mathfrak{A}+\sqrt{m}\mathfrak{B}+\sqrt{n}\mathfrak{C}=0$ is really equivalent to the apparently more general form $\sqrt{l}\mathfrak{A}^\circ+\sqrt{m}\mathfrak{B}^\circ+\sqrt{n}\mathfrak{C}^\circ=0$. In fact, this last is in general a bicircular quartic, and, in regard to it, the before-mentioned theorem of the variable zomal becomes Mr Casey's theorem, that ``the bicircular quartic (and, as a particular case thereof, the circular cubic) is the envelope of a variable circle having its centre on a given conic and cutting at right angles a given circle.'' This theorem is a sufficient basis for the complete theory of the trizomal curve $\sqrt{l}\mathfrak{A}^\circ+\sqrt{m}\mathfrak{B}^\circ+\sqrt{n}\mathfrak{C}^\circ=0$; and it is thereby very easily seen that the curve $\sqrt{l}\mathfrak{A}^\circ+\sqrt{m}\mathfrak{B}^\circ+\sqrt{n}\mathfrak{C}^\circ=0$ can be represented by an equation $\sqrt{l'}\mathfrak{A}'+\sqrt{m'}\mathfrak{B}'+\sqrt{n'}\mathfrak{C}'=0$. But for $\nu>3$ this is not so, and the curve $\sqrt{l}\mathfrak{A}+\&c.=0$ is only a particular form of the curve $\sqrt{l}\mathfrak{A}^\circ+\&c.=0$; and the discussion of this general form is scarcely more difficult than that of the special form $\sqrt{l}\mathfrak{A}+\&c.=0$, included therein.

\setcounter{page}{471}

The investigations in relation to the theory of foci, and to the trizomal and tetrazomal curves where the zomals are circles, are in fact the principal subject of the Memoir, which is divided into Parts as follows: Part I., On Polyzomal Curves in general; Part II., On the Depression of Order, and the Singularities of Polyzomal Curves, with Investigations; Part III., On the Theory of Foci; and Part IV., On the Trizomal and Tetrazomal Curves where the zomals are circles. There is, however, some necessary intermixture of the theories treated of, and the arrangement will appear more in detail from the headings of the several articles.

The paragraphs are numbered continuously through the Memoir. There are four Annexes, relating to questions which it seemed to me more convenient to treat of thus separately.

It is right that I should explain the very great extent to which, in the composition of the present Memoir, I am indebted to Mr Casey's researches. His Paper ``On the Equations and Properties (1) of the System of Circles touching three circles in a plane; (2) of the System of Spheres touching four spheres in space; (3) of the System of Circles touching three circles on a sphere; (4) on the System of Conics inscribed in a conic and touching three inscribed conics in a plane,'' was read to the Royal Irish Academy, April 9, 1866, and is published in their ``Proceedings.'' The fundamental theorem for the equation of the pairs of circles touching three given circles was, previous to the publication of the paper, mentioned to me by Dr Salmon, and I communicated it to Professor Cremona, suggesting to him the problem solved in his letter of March 3, 1866, as mentioned in my paper, ``Investigations in connexion with Casey's Equation,'' Quarterly Math. Journ. vol. viii. 1867, pp.~334---341, [395], and as also appears, Annex No.~IV of the present Memoir.

In connexion with this theorem, I communicated to Mr Casey, in March or April 1867, the theorem No.~164 of the present Memoir, that for any three given circles, centres $A$, $B$, $C$, the equation $BC^2\sqrt{\mathfrak{A}}+CA^2\sqrt{\mathfrak{B}}+AB^2\sqrt{\mathfrak{C}}=0$ (where $BC$, $CA$, $AB$, denote the mutual distances of the points $A$, $B$, $C$) belongs to a Cartesian. Mr Casey, in a letter to me dated 30th April, 1867, informed me of his own mode of viewing the question as follows:---``The general equation of the second order $(a,b,c,f,g,h\!\hbox{\hspace{-0.5em}}\raisebox{0.3ex}{$\scriptstyle\frown$}\!\alpha,\beta,\gamma)^2=0$, where $\alpha$, $\beta$, $\gamma$ are circles, is a bicircular quartic. If we take the equation $(a,b,c,f,g,h\!\hbox{\hspace{-0.5em}}\raisebox{0.3ex}{$\scriptstyle\frown$}\!\lambda,\mu,\nu)^2=0$ in tangential coordinates (that is, when $\lambda$, $\mu$, $\nu$ are perpendiculars let fall from the centres of $\alpha$, $\beta$, $\gamma$ on any line), it denotes a conic; denoting this conic by $F$, and the circle which cuts $\alpha$, $\beta$, $\gamma$ orthogonally by $J$, I proved that, if a variable circle moves with its centre on $F$, and if it cuts $J$ orthogonally, its envelope will be the bicircular quartic whose equation is that written down above;'' and among other consequences, he mentions that the foci of $F$ are the double foci of the quartic, and the points in which $J$ cuts $F$ single foci of the quartic, and also the theorem which I had sent him as to the Cartesian, and he refers to his Memoir on Bicircular Quartics as then nearly finished.

An Abstract of the Memoir as read before the Royal Irish Academy, 10th February, 1867, and published in their Proceedings.

\setcounter{page}{472}

\subsection*{Part I.\ (Nos.\ 1 to 55).---On Polyzomal Curves in General.}

\subsubsection*{Article Nos.\ 1 to 4.\ Definition and Preliminary Remarks.}

\noindent\textbf{1.} As already mentioned, $U$, $V$, \&c.\ denote rational and integral functions $({*}\!\hbox{\hspace{-0.5em}}\raisebox{0.3ex}{$\scriptstyle\frown$}\!x,y,z)^r$, all of the same degree $r$ in the coordinates $(x,y,z)$, and the equation $\sqrt{lU}+\sqrt{mV}+\&c.=0$ then belongs to a polyzomal curve, viz., if the number of the zomes $\sqrt{lU}$, $\sqrt{mV}$, \&c.\ is $=\nu$, then we have a $\nu$-zomal curve.

The radicals, or any of them, may contain rational factors, or be of the form $P\sqrt{Q}$; but in speaking of the curve as a $\nu$-zomal, it is assumed that any two terms, such as $P\sqrt{Q}+P'\sqrt{Q}$, involving the same radical $\sqrt{Q}$, are united into a single term, so that the number of distinct radicals is always $=\nu$; in particular ($r$ being even), it is assumed that there is only one rational term $P$. But the ordinary case, and that which is almost exclusively attended to, is that in which the radicals $\sqrt{U}$, $\sqrt{V}$, \&c.\ are distinct irreducible radicals without rational factors.

\noindent\textbf{2.} The curves $U=0$, $V=0$, \&c.\ are said to be the zomal curves, or simply the zomals of the polyzomal curve $\sqrt{lU}+\sqrt{mV}+\&c.=0$; more strictly, the term zomal would be applied to the functions $U$, $V$, \&c. It is to be noticed, that although the form $\sqrt{U}+\sqrt{V}+\&c.=0$ is equally general with the form $\sqrt{lU}+\sqrt{mV}+\&c.=0$ (in fact, in the former case, the functions $U$, $V$, \&c.\ are considered as implicitly containing the constant factors $l$, $m$, \&c., which are expressed in the latter case), yet it is frequently convenient to express these factors, and thus write the equation in the form $\sqrt{lU}+\sqrt{mV}+\&c.$

For instance, in speaking of any given curves $U=0$, $V=0$, \&c., we are apt, disregarding the constant factors which they may involve, to consider $U$, $V$, \&c.\ as given functions; but in this case the general equation of the polyzomal with the zomals $U=0$, $V=0$, \&c., is of course $\sqrt{lU}+\sqrt{mV}+\&c.=0$.

\noindent\textbf{3.} Anticipating in regard to the cases $\nu=1$, $\nu=2$, the remark which will be presently made in regard to the $\nu$-zomal, that $\sqrt{U}+\sqrt{V}+\&c.=0$ is the curve represented by the rationalised form of this equation, the monozomal curve $\sqrt{U}=0$ is merely the curve $U=0$, viz., this is any curve whatever $U=0$ of the order $r$; and similarly, the bizomal curve $\sqrt{U}+\sqrt{V}=0$ is merely the curve $U-V=0$, viz.\ this is any curve whatever $U=0$, of the order $r$; the zomal curves $U=0$, $V=0$, taken separately, are not curves standing in any special relation to the curve in question $=0$, but $U=0$ may be any curve whatever of the order $r$, and then $V=0$ is a curve of the same order $r$, in involution with the two curves $U=0$, $U=0$; we may, in fact, write the equation $U=0$ under the bizomal form $\sqrt{U}+\sqrt{U}+U=0$. In the case $r$ even, we may, however, notice the bizomal curve $P+\sqrt{U}=0$ ($P$ a rational function of the degree $\frac{1}{2}r$); the rational equation is here $U=(U-P^2)=0$, that is $U=P^2$, viz., $P$ is any curve whatever of the order $\frac{1}{2}r$, and $U=0$ is a curve of the order $r$, touching the given curve $U=0$ at each of its $\frac{1}{2}r^2$ intersections.

\setcounter{page}{474}

It is to be remarked that the order of the $\nu$-zomal curve $\sqrt{U}+\&c.=0$ is $=2^{\nu-1}r$; this is right in the case of the bizomal curve $\sqrt{U}+\sqrt{V}=0$, the order being $=r$, but it fails for the monozomal curve $\sqrt{U}=0$, the order being in this case $r$, instead of $\frac{1}{2}r$, as given by the formula. The two unimportant and somewhat exceptional cases $\nu=1$, $\nu=2$, are thus disposed of, and in all that follows (except in so far as this is in fact applicable to the cases just referred to), $\nu$ may be taken to be $=3$ at least.

\noindent\textbf{4.} It is to be throughout understood that by the curve $\sqrt{lU}+\sqrt{mV}+\&c.=0$ is meant the curve represented by the rationalised equation $\Norm(\sqrt{lU}+\sqrt{mV}+\&c.)=0$, viz.\ the Norm is obtained by attributing to all but one of the zomes $\sqrt{lU}$, $\sqrt{mV}$, \&c., each of the two signs $+$, $-$, and multiplying together the several resulting values of the polyzome; in the case of a $\nu$-zomal curve, the number of factors is thus $=2^{\nu-1}$ (whence, as each factor is of the degree $\frac{1}{2}r$, the order of the curve is $2^{\nu-1}\cdot\frac{1}{2}r$, $=2^{\nu-2}r$, as mentioned above). I expressly mention that, as regards the polyzomal curve, we are not in any wise concerned with the signs of the radicals, which signs are and remain essentially indeterminate; the equation $\sqrt{lU}+\sqrt{mV}+\&c.=0$, is a mere symbol for the rationalised equation, $\Norm(\sqrt{lU}+\sqrt{mV}+\&c.)=0$.

\subsubsection*{Article Nos.\ 5 to 12.\ The Branches of a Polyzomal Curve.}

\setcounter{page}{475}

\noindent\textbf{5.} But we may in a different point of view attend to the signs of the radicals; if for all values of the coordinates we take the symbol $\sqrt{U}$ to signify the positive square root, and consider $\sqrt{U}$, $\sqrt{V}$, \&c.\ as signifying determinately, say the positive values of $\sqrt{U}$, $\sqrt{V}$, \&c.; then each of the several equations $\pm\sqrt{U}\pm\sqrt{V}+\&c.=0$, or, fixing at pleasure one of the signs, suppose that prefixed to $\sqrt{U}$, then each of the several equations $\sqrt{U}\pm\sqrt{V}\pm\&c.=0$, will belong to a branch of the polyzomal curve: a $\nu$-zomal curve has thus $2^{\nu-1}$ branches corresponding to the $2^{\nu-1}$ values respectively of the polyzome. The separation of the branches depends on the precise fixation of the significations of $\sqrt{U}$, $\sqrt{V}$, \&c., and in regard hereto some further explanation is necessary.

\noindent\textbf{6.} When $U$ is real and positive, $\sqrt{U}$ may be taken to be, in the ordinary sense, the positive value of $\sqrt{U}$, and so when $U$ is real and negative, $\sqrt{U}$ may be taken to be $=i$ into the positive value of $\sqrt{-U}$; and the like as regards $\sqrt{V}$, \&c.\ The functions $U$, $V$, \&c.\ are assumed to be real functions of the coordinates; hence, for any real values of the coordinates, $U$, $V$, \&c.\ are real positive or negative quantities, and the significations of $\sqrt{U}$, $\sqrt{V}$, \&c.\ are completely determined.

\noindent\textbf{7.} But the coordinates may be imaginary. In this case the functions $U$, $V$, \&c.\ will for any given values of the coordinates acquire each of them a determinate, in general imaginary, value. If for all real values whatever of $\alpha$, $\beta$, we select once for all one of the two opposite values of $\sqrt{\alpha+\beta i}$, calling it the positive value, and representing it by $\sqrt{\alpha+\beta i}$, then, for any particular values of the coordinates, $U$ being $=\alpha+\beta i$, the value of $\sqrt{U}$ may be taken to be $=\sqrt{\alpha+\beta i}$; and the like as regards $\sqrt{V}$, \&c.\ $\sqrt{U}$, $\sqrt{V}$, \&c.\ have thus each of them a determinate signification for any values whatever, real or imaginary, of the coordinates.

The coordinates of a given point on the curve $\sqrt{lU}+\sqrt{mV}+\&c.=0$, will in general satisfy only one of the equations $\sqrt{U}\pm\sqrt{V}\pm\&c.=0$; that is, the point will belong to one (but in general only one) of the $2^{\nu-1}$ branches of the curve; the entire series of points the coordinates of which satisfy any one of the $2^{\nu-1}$ equations, will constitute the branch corresponding to that equation.

\noindent\textbf{8.} The signification to be attached to the expression $\sqrt{\alpha+\beta i}$ should agree with that previously attached to the like symbol in the case of a positive or negative real quantity; and it should, as far as possible, be subject to the condition of continuity, viz., as $\alpha+\beta i$ passes continuously to $\alpha'+\beta' i$, so $\sqrt{\alpha+\beta i}$ should pass continuously to $\sqrt{\alpha'+\beta' i}$; but (as is known) it is not possible to satisfy universally this condition of continuity; viz., if for facility of explanation we consider $(\alpha,\beta)$ as the coordinates of a point in a plane, and imagine this point to describe a closed curve surrounding the origin or point $(0,0)$, then it is not possible so to define $\sqrt{\alpha+\beta i}$ that this quantity, varying continuously as the point moves along the curve, shall, when the point has made a complete circuit, resume its original value. The signification to be attached to $\sqrt{\alpha+\beta i}$ is thus in some measure arbitrary, and it would appear that the division of the curve into branches is affected by a corresponding arbitrariness, but this arbitrariness relates only to the imaginary branches of the curve: the notion of a real branch is perfectly definite.

\noindent\textbf{9.} It would seem that a branch may be impossible for any series whatever of points real or imaginary. Thus, in the bizomal curve $\sqrt{U}+\sqrt{V}=0$, the branch $\sqrt{U}+\sqrt{V}=0$ is impossible. In fact, for any point whatever, real or imaginary, of the curve, we have $U=V$, and therefore $\sqrt{U}=\sqrt{V}$; the point thus belongs to the other branch $\sqrt{U}-\sqrt{V}=0$, not to the branch $\sqrt{U}+\sqrt{V}=0$; the only points belonging to the last-mentioned branch are the isolated points for which simultaneously $\sqrt{U}=0$, $\sqrt{V}=0$; viz., the points of intersection of the two curves $U=0$, $V=0$.

\noindent\textbf{10.} It is not clear to me whether the case is the same in regard to the branch $\sqrt{U}+\sqrt{V}+\sqrt{W}=0$ of a trizomal curve. In fact, for each point of the curve $\sqrt{U}+\sqrt{V}+\sqrt{W}=0$ we have $(U-V-W)^2=4VW$, and therefore, $U-V-W=\pm 2\sqrt{V}\sqrt{W}$; there may very well be points for which the sign is $+$; that is, points for which $U=V+W+2\sqrt{V}\sqrt{W}$, and for these points we have $\pm\sqrt{U}=\sqrt{V}+\sqrt{W}$; for real values of the coordinates the sign on the left hand must be $+$ (for otherwise the two sides would have opposite signs), but there is no apparent reason, or at least no obviously apparent reason, why this should be so for imaginary values of the coordinates, and if the sign be in fact $-$, then the point will belong to the branch $\sqrt{U}+\sqrt{V}+\sqrt{W}=0$.

\noindent\textbf{11.} But the branch in question is clearly impossible for any series of real points; so that, leaving it an open question whether the epithet ``impossible'' is to be understood to mean impossible for any series of real points (that is, as a mere synonym of imaginary), or whether it is to mean impossible for any series whatever of points real or imaginary, I say that the branch $\sqrt{U}+\sqrt{V}+\sqrt{W}=0$ of a trizomal curve is an impossible branch.

\setcounter{page}{476}

\noindent\textbf{12.} Consider a $\nu$-zomal curve $\sqrt{lU}+\sqrt{mV}+\&c.=0$; the equations of the several branches are the equations $\sqrt{U}\pm\sqrt{V}\pm\&c.=0$, obtained by attributing to the radicals $\sqrt{V}$, $\sqrt{W}$, \&c.\ the signs $+$ or $-$ at pleasure, the sign of $\sqrt{U}$ being fixed. If in the equation of any branch we write $\sqrt{V}$, \&c.\ for those radicals which have the sign $-$, then the foregoing equations break up into the more simple equations $\sqrt{U}+\&c.=0$, $\sqrt{V}+\&c.=0$, which are the equations of certain branches of the curves $\sqrt{lU}+\&c.=0$, $\sqrt{mV}+\&c.=0$, respectively, and conversely each of the intersections of these two curves is a point situate simultaneously on some two branches of the original $\nu$-zomal curve $\sqrt{lU}+\sqrt{mV}+\&c.=0$.

Hence, partitioning in any manner the $\nu$-zome $\sqrt{lU}+\sqrt{mV}+\&c.$ into an $\alpha$-zome, $\sqrt{lU}+\&c.$ and a $\beta$-zome $\sqrt{mV}+\&c.$ ($\alpha+\beta=\nu$), and writing down the equations $\sqrt{lU}+\&c.=0$, $\sqrt{mV}+\&c.=0$ of an $\alpha$-zomal curve and a $\beta$-zomal curve respectively, each of the intersections of these two curves is a point situate simultaneously on two branches of the $\nu$-zomal curve; and the points situate simultaneously on two branches of the $\nu$-zomal curve are the points of intersection of the several pairs of an $\alpha$-zomal curve and a $\beta$-zomal curve, which can be formed by any bipartition of the $\nu$-zome.

\noindent\textbf{13.} There are two cases to be considered:---First, when the parts are $1$, $\nu-1$ ($\nu-1$ is $>1$, except in the case $\nu=2$, which may be excluded from consideration), or say when the $\nu$-zome is partitioned into a zome and antizome. Secondly, when the parts $\alpha$, $\beta$, are each $>1$ (this implies $\nu=4$ at least), or say when the $\nu$-zome is partitioned into a pair of complementary parazomes.

\noindent\textbf{14.} To fix the ideas, take the tetrazomal curve $\sqrt{U}+\sqrt{V}+\sqrt{W}+\sqrt{T}=0$, and consider first a point for which $\sqrt{U}=0$, $\sqrt{V}+\sqrt{W}+\sqrt{T}=0$. The Norm is the product of ($2^3=$) 8 factors; selecting hereout the factors $\sqrt{U}+\sqrt{V}+\sqrt{W}+\sqrt{T}$, $\sqrt{U}-\sqrt{V}-\sqrt{W}-\sqrt{T}$, let the product of these be called $F$, and the product of the remaining six factors be called $G$; the rationalised equation of the curve is therefore $FG=0$. The derived equation is $GdF+FdG=0$; at the point in question $\sqrt{U}=0$, $\sqrt{V}+\sqrt{W}+\sqrt{T}=0$; $G$ and $dG$ are each of them finite (that is, they neither vanish nor become infinite), but we have $F=0$, $dF=dU\cdot(\sqrt{V}+\sqrt{W}+\sqrt{T})=0$, that is, the derived equation becomes $0=0$, or the point in question is an ordinary point on the curve.

\setcounter{page}{478}

\noindent\textbf{15.} Consider, secondly, a point for which $\sqrt{U}+\sqrt{V}=0$, $\sqrt{W}+\sqrt{T}=0$; to form the Norm, taking in this case the two factors $\sqrt{U}+\sqrt{V}+\sqrt{W}+\sqrt{T}$, $\sqrt{U}+\sqrt{V}-\sqrt{W}-\sqrt{T}$, let their product $=(\sqrt{U}+\sqrt{V})^2-(\sqrt{W}+\sqrt{T})^2$ be called $F$, and the product of the remaining six factors be called $G$; the rationalised equation is $FG=0$, and the derived equation is $FdG+GdF=0$. At the point in question $G$ and $dG$ are each of them finite (that is, they neither vanish nor become infinite), but we have $F=0$, $dF=(\sqrt{U}+\sqrt{V})(dU/\sqrt{U}+dV/\sqrt{V})-(\sqrt{W}+\sqrt{T})(dW/\sqrt{W}+dT/\sqrt{T})=0$, that is, the derived equation becomes identically $=0$; the point in question is thus a singular point, and it is easy to see that it is in fact a node, or ordinary double point, on the tetrazomal curve. And similarly, each of the points of intersection of the two curves $\sqrt{U}+\sqrt{V}=0$, $\sqrt{W}+\sqrt{T}=0$ is a node on the tetrazomal curve.

\noindent\textbf{16.} The proofs in the foregoing two examples respectively are quite general, and we may, in regard to a $\nu$-zomal curve, enunciate the results as follows, viz., in a $\nu$-zomal curve, the points situate simultaneously on two branches are either the intersections of a zomal curve and its antizomal curve, or else they are the intersections of a pair of complementary parazomal curves. In the former case, the points in question are ordinary points on the $\nu$-zomal, but they are points of contact of the $\nu$-zomal with the zomal; it may be added, that the intersections of the zomal and antizomal, each reckoned twice, are all the intersections of the $\nu$-zomal and zomal. In the latter case, the points in question are nodes of the $\nu$-zomal; it may be added, that the $\nu$-zomal has not, in general, any nodes other than the points which are thus the intersections of a pair of complementary parazomals, and that it has not in general any cusps.

\subsubsection*{Article Nos.\ 17 to 21.\ Singularities of a $\nu$-zomal Curve.}

\setcounter{page}{479}

\noindent\textbf{17.} It has been already shown that the order of the $\nu$-zomal curve is $=2^{\nu-2}r$. Considering the case where $\nu$ is $=3$ at least, the curve has been represented as the partial intersection of two surfaces, the order of the complete intersection being $=2^{\nu-1}\cdot\frac{1}{2}r$, $=2^{\nu-2}r$; and the curve has thus a deficiency equal to the number of intersections which are wanting to make up the complete intersection; or the deficiency is $=2^{\nu-2}r(2^{\nu-2}r-1)-2\cdot 2^{\nu-3}r(2^{\nu-3}r-1)$, $=2^{2\nu-5}r^2$.

\noindent\textbf{18.} Considering first a zome and its antizome, say the zome $U=0$ and the antizome $\sqrt{V}+\sqrt{W}+\&c.=0$ (a $(\nu-1)$-zomal curve); the zome is of the order $r$, and the antizome of the order $2^{\nu-3}r$; they intersect therefore in $2^{\nu-3}r^2$ points, each of which is a point of contact of the $\nu$-zomal with the zome. The number of contacts with all the zomes is thus $=\nu\cdot 2^{\nu-3}r^2$.

\noindent\textbf{19.} Considering next a pair of complementary parazomal curves, an $\alpha$-zomal and a $\beta$-zomal respectively ($\alpha+\beta=\nu$), these are of the orders $2^{\alpha-2}r$ and $2^{\beta-2}r$ respectively, and they intersect therefore in $2^{\alpha+\beta-4}r^2=2^{\nu-4}r^2$ points, nodes of the $\nu$-zomal. This number is independent of the particular partition $(\alpha,\beta)$, and the $\nu$-zomal has thus this same number, $2^{\nu-4}r^2$, of nodes in respect of each pair of complementary parazomals; hence the total number of nodes is $=2^{\nu-4}r^2$ into the number of pairs of complementary parazomals.

For the partition $(\alpha,\beta)$ the number of pairs is $=[\nu]^\alpha\div[\alpha]^\alpha[\beta]^\beta$, or when $\alpha=\beta$, which of course implies $\nu$ even, it is one-half of this; extending the summation from $\alpha=2$ to $\alpha=\nu-2$, each pair is obtained twice, and the number of pairs is thus $=\frac{1}{2}\{(1+1)^\nu-\nu-\nu-1-1\}$; the sum extended from $\alpha=0$ to $\alpha=\nu$ is $(1+1)^\nu=2^\nu$, but we thus include the terms $1$, $\nu$, $\nu$, $1$, which are together $=2\nu+2$, hence the correct value of the sum is $=2^\nu-2\nu-2$, and the number of pairs is the half of this $=2^{\nu-1}-\nu-1$.

Hence the number of nodes of the $\nu$-zomal curve is $=(2^{\nu-1}-\nu-1)2^{\nu-4}r^2$.

\setcounter{page}{480}

\noindent\textbf{20.} The $\nu$-zomal is thus a curve of the order $2^{\nu-2}r$, with $(2^{\nu-1}-\nu-1)2^{\nu-4}r^2$ nodes, but without cusps; the class is therefore $=2^{\nu-2}r[(\nu+1)r-2]$, and the deficiency is $=2^{\nu-4}r[(\nu+1)r-6]+1$.

These are the general expressions, but even when the zomal curves $U=0$, $V=0$, \&c., are given, then writing the equation of the $\nu$-zomal under the form $\sqrt{lU}+\sqrt{mV}+\&c.=0$, the constants $l:m:\&c.$, may be so determined as to give rise to nodes or cusps which do not occur in the general case; the formulae will also undergo modification in the particular cases next referred to.

\subsubsection*{Article Nos.\ 22 to 27.\ Special Case where all the Zomals have a Common Point or Points.}

\noindent\textbf{21.} Consider the case where the zomals $U=0$, $V=0$ have all of them any number, say $k$, of common intersections---these may be referred to simply as the common points. Each common point, being a $2^{\alpha-2}$-tuple point on the $\alpha$-zomal and a $2^{\beta-2}$-tuple point on the $\beta$-zomal, counts as $2^{\alpha+\beta-4}$, $=2^{\nu-4}$ intersections, and the $k$ common points count as $2^{\nu-4}k$ intersections; the number of the remaining intersections is therefore $=2^{\nu-4}(r^2-k)$, each of which is a node on the $\nu$-zomal curve; and we have thus in all $2^{\nu-4}(2^{\nu-1}-\nu-1)(r^2-k)$ nodes.

\noindent\textbf{22.} There are, besides, the $k$ common points, each of them a $2^{\nu-2}$-tuple point on the $\nu$-zomal, and therefore each reckoning as $\frac{1}{2}2^{\nu-2}(2^{\nu-2}-1)$, $=2^{2\nu-5}-2^{\nu-3}$ double points, or together as $(2^{2\nu-5}-2^{\nu-3})k$ double points.

Reserving the term node for the above-mentioned nodes or proper double points, and considering, therefore, the double points (dps.) as made up of the nodes and of the $2^{\nu-2}$-tuple points, the total number of dps.\ is thus
\begin{align*}
&2^{\nu-4}(2^{\nu-1}-\nu-1)(r^2-k)+(2^{2\nu-5}-2^{\nu-3})k,\\
&\qquad=2^{\nu-4}(2^{\nu-1}-\nu-1)r^2+[(\nu+1)2^{\nu-4}-2^{\nu-3}]k;
\end{align*}
or finally this is $=2^{\nu-4}\{(2^{\nu-1}-\nu-1)r^2+(\nu-1)k\}$; so that there is a gain $=2^{\nu-4}(\nu-1)k$ in the number of dps.\ arising from the $k$ common points.

There is, of course, in the class a diminution equal to twice this number, or $2^{\nu-3}(\nu-1)k$; and in the deficiency a diminution equal to this number, or $2^{\nu-4}(\nu-1)k$.

\noindent\textbf{23.} The zomal curves $U=0$, $V=0$, \&c., may all of them pass through the same $r^2$ points; we have then $k=r^2$, and the expression for the number of dps.\ is $=(2^{2\nu-5}-2^{\nu-3})r^2$, viz., this is $=2^{\nu-3}(2^{\nu-2}-1)r^2$.

But in this case the dps.\ are nothing else than the $r^2$ common points, each of them a $2^{\nu-2}$-tuple point, the $\nu$-zomal curve in fact breaking up into a system of $2^{\nu-2}$ curves of the order $r$, each passing through the $r^2$ common points.

This is easily verified, for if $\Theta=0$, $\Phi=0$ are some two curves of the order $r$, then, in the present case, the zomal curves are curves in involution with these curves; that is, they are curves of the form $l\Theta+l'\Phi=0$, and the $\nu$-zomal curve $\sqrt{l(l\Theta+l'\Phi)}+\&c.=0$ is thus satisfied by $l\Theta+l'\Phi=0$, that is, it breaks up into the system of curves $l\Theta+l'\Phi=0$, where $l:l'$ has $2^{\nu-2}$ values.

\setcounter{page}{481}

\noindent\textbf{24.} A more important case is when the zomal curves are each of them in involution with the same two given curves, one of them of the order $r$, the other of an inferior order. Let $\Theta=0$ be a curve of the order $r$, $\Phi=0$ a curve of an inferior order $r-s$; $L=0$, $M=0$, \&c., curves of the order $s$; then the case in question is when the zomal curves are of the form $\Theta+L\Phi=0$, $\Theta+M\Phi=0$, \&c., the equation of the $\nu$-zomal is
\[
\sqrt{l(\Theta+L\Phi)}+\sqrt{m(\Theta+M\Phi)}+\&c.=0,
\]
where $l$, $m$, \&c.\ are constants. This is the most convenient form for the equation, and by considering the functions $L$, $M$, \&c.\ as containing implicitly the factors $l^{-\frac{1}{2}}$, $m^{-\frac{1}{2}}$, \&c.\ respectively, we may take it to include the form
\[
\sqrt{\Theta+L\Phi}+\sqrt{\Theta+M\Phi}+\&c.=0,
\]
which last has the advantage of being immediately applicable to the case where any one or more of the constants $l$, $m$, \&c.\ may be $=0$.

\noindent\textbf{25.} In the case now under consideration we have the $r(r-s)$ points of intersection of the curves $\Theta=0$, $\Phi=0$ as common points of all the zomals. Hence, putting in the foregoing formula $k=r(r-s)$, we have a $\nu$-zomal curve of the order $2^{\nu-2}r$, having with each zomal $2^{\nu-2}rs$ contacts, or with all the zomals $2^{\nu-2}rs\nu$ contacts; the number of nodes is $=2^{\nu-4}(2^{\nu-1}-\nu-1)\{r^2-r(r-s)\}$, $=2^{\nu-4}(2^{\nu-1}-\nu-1)rs$; and there is a diminution in the class $=2^{\nu-3}(\nu-1)r(r-s)$, and in the deficiency $=2^{\nu-4}(\nu-1)r(r-s)$.

\end{document}
